%%%%%%%%%%%%%%%%%%%%%%%%%%%%%%%%%%%%%%%%%%%%%%%%%%%%%%%%%%%%%%%%%%%%%%%%
%                                                                      %
%                               Chapter 11                             %
%                                                                      %
%%%%%%%%%%%%%%%%%%%%%%%%%%%%%%%%%%%%%%%%%%%%%%%%%%%%%%%%%%%%%%%%%%%%%%%%

\chapter{Techniques of Integration}

Chapter 6 introduced the integral.  There it was defined numerically,
as the limit of approximating Riemann sums.  Evaluating integrals by
applying this basic definition tends to take a long time if a high
level of accuracy is desired.  If one is going to evaluate integrals
at all frequently, it is thus important to find {\bf techniques of
  integration} for doing this efficiently.  For instance, if we
evaluate a function at the midpoints of the subintervals, we get much
faster convergence than if we use either the right or left endpoints
of the subintervals.

A powerful class of techniques is based on the observation made at the
end of chapter 6, where we saw that the fundamental theorem of
calculus gives us a second way to find an integral, using
antiderivatives. While a Riemann sum will usually give us only an
approximation to the value of an integral, an antiderivative will give
us the exact value.  The drawback is that antiderivatives often can't
be expressed in {\bf closed form}\index{closed form
  solution}\index{solution!closed form}---that is, as a {\bf
  formula} in terms of named functions.  Even when antiderivatives can
be so expressed, the formulas are often difficult to find.
Nevertheless, such a formula can be so powerful, both computationally
and analytically, that it is often worth the effort needed to find it.
In this chapter we will explore several techniques for finding the
antiderivative of a function given by a formula.

We will conclude the chapter by developing a numerical
method---Simpson's rule---that gives a good estimate for the value of
an integral with relatively little computation.

\newpage
%%%%%%%%%%%%%%%%%%%%%%%%%%%%%%%%%%%%%%%%%%%%%%%%%%%%%%%%%%%%%%%%%%%%%%%%
%                                                                      %
%                               Section 1                              %
%                                                                      %
%%%%%%%%%%%%%%%%%%%%%%%%%%%%%%%%%%%%%%%%%%%%%%%%%%%%%%%%%%%%%%%%%%%%%%%%

\section{Antiderivatives}

\subsection{Definition}

Recall that we say $F$ is an {\bf
  antiderivative}\index{antiderivative!definition of} of $f$
if $F' = f$.  Here are some examples.
\begin{center}
\begin{tabular}{rcccccccccc}
	{\sc function}: && $x^2$ && $1/y$ && $\sin{u}$ && 
		$2\sin{t} \, \cos{t}$ && $ 2^z$ \\ [3pt]
	&&$\updownarrow$ && $\updownarrow$ && $\updownarrow$ && 
		$\updownarrow$ && $\updownarrow$  \\ [4pt]
	{\sc antiderivative}:&&$\dfrac{x^3}{3}$&& 
		$\ln{y}$ && $-\cos{u}$ && $ \sin^2{t}$ && 
		$\dfrac{2^z}{\ln{2}}$
\end{tabular}
\end{center}
Notice that you go up ($\uparrow$) from the bottom row 
%
\mar{Undo a differentiation}
%
to the top by carrying out a differentiation\index{differentiation}.
To go down ($\downarrow$) you must ``undo'' that differentiation.  The
process of reversing, or undoing, a differentiation has come to be
called {\bf antidifferentiation}\index{antidifferentiation}.  You
should differentiate each function on the bottom row to check that it
is an antiderivative of the function above it.


While a function can have only one derivative, 
%
\mar{A function has\\ many antiderivatives}	
%
it has many antiderivatives\index{antiderivative}.  For example,
$1-\cos{u}$ and $99 - \cos{u}$ are also antiderivatives of the
function $\sin{u}$ because
	$$
	(1-\cos{u})' = \sin{u} = (99 - \cos{u})'.
	$$
In fact, every function $C - \cos{u}$ is an antiderivative of
$\sin{u}$, for any constant $C$ whatsoever.  This observation is true
in general.  That is, if $F$ is an antiderivative of a function $f$,
then so is $F + C$, for any constant $C$.  This follows from the
addition rule for derivatives\index{derivative!addition rule}:
	$$
	(F + C)' = F' + C' = F' + 0 = f.
	$$

It is tempting to claim 
%
\mar{A caution}
%
the converse---that {\em every\/} antiderivative of $f$ is equal to $F
+ C$, for some appropriately chosen value of $C$.  In fact, you will
often see this statement written.  The statement is true, though, only
for continuous functions.  If the function $f$ has breaks in its
domain, then there will be more antiderivatives than those of the form
$F+C$ for a {\em single} constant~$C$---over each piece of the domain
of $f$, $F$ can be modified by a {\em different} constant and still
yield an antiderivative for $f$.  Exercises 18, 19, and 20 at the end
of this section explore this for a couple of cases.  If $f$ is
continuous, though, $F + C$ will cover all the possibilities, and
%
\mar{What the `$+\ C$' term really means} 
%
we sometimes say that $F + C$ is {\em the\/} antiderivative of $f$.
For the sake of keeping a compact notation, we will even write this
when the domain of $f$ consists of more than one interval.  You should
understand, though, that in such cases, over each piece $F$ can be
modified by a different constant

For future reference we collect a list 
%
\mar{Antiderivatives\\ of basic functions}	
%
of basic functions whose antiderivatives\index{antiderivative} we
already know.  Remember that each antiderivative in the table can have
an arbitrary constant added to it.
	$$
	\begin{array}{ccc}
	\mbox{function} & \qquad &	\mbox{antiderivative} \\[2pt]
	\hline \\ [-9pt]
	\begin{array}{l}
		x^p 	\\ [12pt]
		1/x 	\\ [2pt]
		\sin{x} \\ [2pt]
		\cos{x} \\ [2pt]
		e^x	\\ [6pt]
		b^x	
	\end{array} 
	& & 
	\begin{array}{l} 
		\hphantom{-}\dfrac{x^{p+1}}{p+1}, \quad p \neq -1 \\[10pt]
		\hphantom{-}\ln{x}\\ [2pt]
			-\cos{x}\\ [2pt]
		\hphantom{-}\sin{x}\\ [2pt]
		\hphantom{-}e^x \\[2pt]
		\hphantom{-}\dfrac{b^x}{\ln{b}}
	\end{array}
	\end{array}
	$$

All of these antiderivatives are easily verified and could have been
derived with at most a little trial and error fiddling to get the
right constant.  You should notice one incongruity: the function $1/x$
is defined for all $x\neq 0$, but its listed antiderivative, $\ln{x}$,
is only defined for $x>0$.  In exercise 18 (page~\pageref{log}) you
will see how to find antiderivatives for $1/x$ over its entire domain.

Two of our basic functions---$\ln{x}$ and $\tan{x}$---do not appear in
the left column of the table.  This happens because there is no simple
multiple of a basic function whose derivative is equal to either
$\ln{x}$ or $\tan{x}$.  It turns out that these functions {\em do\/}
have antiderivatives, though, that can be expressed as more
complicated combinations of basic functions.  In fact, by
differentiating $x \ln{x} - x$ you should be able to verify that it is
an antiderivative of $\ln{x}$.  Likewise, $- \ln(\cos{x})$ is an
antiderivative of $\tan{x}$.\label{tangent} It would take a long time
to stumble on these antiderivatives by inspection or by trial and
error.  It is the purpose of later sections to develop techniques
which will enable us to discover antiderivatives like these quickly
and efficiently.  In particular, the antiderivative of $\ln{x}$ is
derived in chapter~11.3 on page~\pageref{lnder}, while the
antiderivative of $\tan{x}$ is derived in chapter~11.5 on
page~\pageref{tander}.

There are a couple of other functions that don't appear in the above
table whose antiderivatives are needed frequently enough that they
should become part of your repertoire of elementary functions that you
recognize immediately:

\newpage

\mbox{}\vspace{-20pt}
$$
	\begin{array}{ccc}
	\mbox{function} & \qquad &	\mbox{antiderivative} \\[2pt]
	\hline \\ [-9pt]
	\begin{array}{l}
		\dfrac{1}{\sqrt{1-x^2}}	\\ [10pt]
		\dfrac{1}{1+x^2}\label{arcsin deriv}	
	\end{array} 
	& & 
	\begin{array}{l} 
		\hphantom{-} \arcsin{x} \\ [10pt]
		\hphantom{-} \arctan{x}  
	\end{array}
	\end{array}
	$$ 
The antiderivatives are inverse trigonometric functions, which we've
had no need for until now.  We introduce them immediately below.  They
are examples of functions that occur more often for their
antiderivative properties than for themselves.  Note that the
derivatives of the inverse trigonometric functions have no obvious
reference to trigonometric relations.  In fact, they often occur in
settings where there are no triangles or periodic functions in sight.
Let's see how the derivatives of these inverse functions are derived.

\subsection{Inverse Functions}

We discussed inverse functions in chapter 4.4.  Here's a quick summary
of the main points made there.  Two functions $f$ and $g$ are {\bf
  inverses}\index{function!inverse} if
\begin{align*}
f(g(a)) &= a,  \\
\text{and} \quad g(f(b)) &= b,
\end{align*}
for every $a$ in the domain\index{domain} of $g$ and every $b$ in the
domain of~$f$.  It follows that the {\em range} of~$f$ is the same as
the {\em domain} of~$g$, and vice versa.  We write $g = f^{-1}$ to
indicate that $g$ is the inverse of~$f$, and $f = g^{-1}$ to indicate
that $f$ is the inverse of~$g$.

The graphs of $y = f(x)$ and $y = g(x)$ are mirror reflections about
the line $y=x$.  As the figure below shows, the mirror image of a
point with coordinates $(a, b)$ is the point with coordinates $(b,
a)$.

\vspace{130pt}

\special{psfile=../ps/calc2/ch11-1.ps}

\newpage

\noindent
This connection between the graphs is a direct translation of the
definitions into graphical language, since
\begin{align*}
\text{$(a, b)$ is on the graph of $y=g(x)$} &\Longleftrightarrow
     g(a) = b \qquad \text{(definition of graph)} \\
  &\Longleftrightarrow  f(b)=a \qquad \text{(definition of inverse}) \\
  &\Longleftrightarrow  \text{$(b, a)$ is on the graph of $y=f(x)$}.
\end{align*}

Because of the connection between the graphs, it follows immediately
that if the graph of $f$ is locally linear at the point $(b, a)$ with
slope $m$, then the graph of $g$ will be locally linear at the point
$(a, b)$ with slope $1/m$.  Algebraically, this is expressed as
\[
g'(a)=1/f'(b),
\]
where $a=f(b)$ and $b=g(a)$.  We get same result by differentiating
the expression $f(g(x))=x$, using the chain rule:
\[
1=x'=(f(g(x))'=f'(g(x))g'(x),
\]
and therefore
\[
g'(x)=\dfrac{1}{f'(g(x))}
\]
for any value of $x$ for which $g(x)$ is defined.  

\subsubsection{Inverse trigonometric functions}

{\bf The arcsine function}\index{function!arcsine}\quad In the
discussion in chapter 4 we saw that a function has an inverse only
when it is {\bf one-to-one}, so if we want an inverse, we often have
to restrict the domain of a function to a region where it is
one-to-one. This is certainly the case with the sine function, which
takes the same value infinitely many times.  The standard choice of
domain on which the sine function is one-to-one is $[-\pi/2, \pi/2]$.
Over this interval the sine function increases from $-1$ to $1$.  We
can then define an inverse function, which we call the {\bf arcsine
  function}, written $\arcsin{x}$, whose domain is the interval
$[-1,1]$, and whose range is $[-\pi/2, \pi/2]$. Since the sine
function is strictly increasing on its domain, the arcsine function
will be strictly increasing on its domain as well---do you see why
this has to be?

Another notation for the arcsine function is $\sin^{-1}{x}$; this form
commonly appears on one of the buttons on a calculator.

\newpage

\mbox{}

\vspace*{140pt}

\special{psfile=../ps/calc2/ch11-2a.ps}
\special{psfile=../ps/calc2/ch11-2b.ps}

To find the derivative of the arcsine function, let $f(x)=\sin{x}$,
and let $g(x)=\arcsin{x}$.  Then by the remarks above, we have
\[
g'(x) = \dfrac{1}{f'(g(x))} = \dfrac{1}{\cos{(g(x))}} 
    = \dfrac{1}{\cos(\arcsin{x})}.
\]
The function $\cos(\arcsin{x})$ can be expressed in another form,
which we will obtain two ways, one algebraic, the other geometric.
Both perspectives are useful.

\noindent
{\bf The algebraic approach.}\quad Recall that for any input $t$,
$\sin^2{t} + \cos^2{t}=1$.  This can be solved for $\cos{t}$ as
$\cos{t}=\pm\sqrt{1-\sin^2{t}}$.  That is, the cosine of anything is
the square root of 1 minus the square of the sine of that input, with
a possible minus sign needed out front, depending on the
context. Since the output of the arcsine function lies in the range
$[-\pi/2, \pi/2]$, and the cosine function is positive (or 0) for
numbers in this interval, it follows that $\cos(\arcsin{x})\geq 0$ for
any value of $x$ in the domain of the arcsine function.  Therefore,
\[
\cos(\arcsin{x}) = \sqrt{1-\sin^2(\arcsin{x})}
  =\sqrt{1-(\sin(\arcsin{x}))^2}=\sqrt{1-x^2},
\]
since $\sin(\arcsin{x})=x$ by definition of inverse functions.  It
follows that
\[
(\arcsin{x})'=\dfrac{1}{\cos{(\arcsin{x})}}=\dfrac{1}{\sqrt{1-x^2}},
\]
as we indicated above on page~\pageref{arcsin deriv}.

\noindent
{\bf The geometric approach}\quad Introduce a new variable $\theta =
\arcsin{x}$, so that $x=\sin{\theta}$.  We can represent these
relationships in the following picture:

\begin{center}
\begin{picture}(180,85)(-10,5)\label{picture}
\bezier{150}(30,10)(29.2,16.89)(26.83,23.42)
\put(165,47){$x$}
\put(78,57){$1$}
\put(80,5){\makebox(0,0)[t]{$\sqrt{1-x^2}$}}
\put(32,15){$\theta$}
\thicklines
\put(0,10){\line(1,0){160}}
\put(160,10){\line(0,1){80}}
\put(0,10){\line(2,1){160}}
%\put(20,10){\line(0,1){2.8}}
%\put(17.88,18.94){\line(1,-2){1.17}}
%\put(20,12.8){\line(-1,4){2.12}}
\end{picture}
\end{center}
Notice that we have labelled the side opposite the angle $\theta$ as
$x$ and the hypotenuse as 1.  This ensures that $\sin\theta=x$.  By
the Pythagorean theorem, the remaining side must then be
$\sqrt{1-x^2}$.  From this picture it is then obvious that
$\cos(\arcsin{x})=\cos\theta=\sqrt{1-x^2}/1=\sqrt{1-x^2},$ as before.


\medskip
\noindent
{\bf The arctangent function}\index{function!arctangent}\quad To get
an inverse for the tangent function, we again need to limit its
domain.  Here the standard choice is to restrict it to the interval
$(-\pi/2, \pi/2)$ ({\em not} including the endpoints this time, since
$\tan(-\pi/2)$ and $\tan(\pi/2)$ aren't defined).  Over this domain
the tangent function increases from $-\infty$ to $+\infty$.  We can
then define the inverse of the tangent, called the {\bf arctangent
  function}, written $\arctan{x}$, whose domain is the interval
$(-\infty, \infty)$, and whose range is $(-\pi/2, \pi/2)$. Again, both
functions are increasing over their domains.

\vspace*{2.8in}
\special{psfile=../ps/calc2/ch11-3a.ps}
\special{psfile=../ps/calc2/ch11-3b.ps}

To find the derivative, we proceed as we did with $\arcsin{x}$,
letting $f(x)=\tan{x}$, and $g(x)=\arctan{x}$.  This time we get
\[
g'(x) = \dfrac{1}{f'(g(x))} = \dfrac{1}{\sec^2{(g(x))}} 
  =\dfrac{1}{\sec^2(\arctan{x})}.
\]
This expression also has a different form; we obtain it from another
trigonometric identity, as we did before, deriving the desired
result algebraically and geometrically.

\noindent
{\bf Algebraic}\quad We start as before with the identity $\sin^2{t} +
\cos^2{t}=1$.  Dividing through by $\cos^2{t}$, we get the equivalent
identity $\tan^2{t}+1=\sec^2{t}$.  That is, the square of the secant
of any input is just 1 plus the square of the tangent applied to the
same input.  In particular,
\[
\sec^2{(\arctan{x})} = \tan^2(\arctan{x}) + 1
    =(\tan(\arctan{x}))^2 + 1 = x^2+1;
\]
since $\tan(\arctan{x}) = x$.  It follows that
\[
(\arctan{x})' = \dfrac{1}{\sec^2(\arctan{x})} = \dfrac{1}{1+x^2},
\]
as we indicated above on page~\pageref{arcsin deriv}.  

\noindent
{\bf Geometric}\quad Let $\theta=\arctan{x}$, so $x=\tan\theta$.
Again we can draw and label a triangle reflecting these relationships:
\begin{center}
\begin{picture}(180,85)(-10,0)\label{picture2}
\bezier{150}(30,10)(29.2,16.89)(26.83,23.42)
\put(165,47){$x$}
\put(78,57){\makebox(0,0)[r]{$\sqrt{1+x^2}$}}
\put(80,5){\makebox(0,0)[t]{$1$}}
\put(32,15){$\theta$}
\thicklines
\put(0,10){\line(1,0){160}}
\put(160,10){\line(0,1){80}}
\put(0,10){\line(2,1){160}}
%\put(20,10){\line(0,1){2.8}}
%\put(17.88,18.94){\line(1,-2){1.17}}
%\put(20,12.8){\line(-1,4){2.12}}
\end{picture}
\end{center}
Notice that this time we have labelled the side opposite the angle
$\theta$ as $x$ and the adjacent side as 1 to ensure that
$\tan\theta=x$.  Again by the Pythagorean theorem, the hypotenuse must
be $1+x^2$.  From this picture it is then obvious that
$\sec(\arctan{x})=\sec\theta=(\sqrt{1+x^2})/1=\sqrt{1+x^2},$ so
$\sec^2(\arctan{x})=1+x^2$ again.


\subsection{Notation}

According to the fundamental theorem of calculus\index{fundamental
  theorem of calculus} (see chapter~6.4), every accumulation
function\index{accumulation function!as antiderivative}
	$$
	A(x) = \int_a^x f(t) \, dt
	$$
is an antiderivative of the function $f$, no matter at what point $t = a$ the accumulation begins---so long as the function is defined over the entire interval from $t=a$ to $t=x$.  (This is an important caution when dealing with functions like $1/x$, for instance, for which the integral from, say, $-1$ to 1 makes no sense.)  In other words, the expression
	$$
	\int_a^x f(t) \, dt
	$$
represents an antiderivative of $f$.  
	\mar{Write an antiderivative as an integral}
The influence of the fundamental theorem is so pervasive that this expression---with the ``limits of integration'' $a$ and $x$ omitted---is used to denote an antiderivative:\index{integral}
\begin{center}
\begin{picture}(310,45)
\put(0,0){\framebox(310,45)
{{\bf Notation}: The antiderivative of $f$ is ${\dint} f(x) \, dx$.
}}
\end{picture}
\end{center}
With this new notation\index{notation!antiderivative} the antiderivatives we have listed so far can be written \index{antiderivative} 
	\mar{\vspace{82pt} The basic antiderivatives again}
in the following form.\label{table}
\begin{align*}
\int x^p \, dx 
   &= \dfrac{1}{p+1} \, x^{p+1} + C  \qquad\mbox{ ( $p\neq -1$)} \\
\int \dfrac{1}{x} \, dx &= \ln{x} + C\\
\int \sin{x} \, dx &= - \cos{x} + C \\
\int \cos{x} \, dx &= \sin{x} + C \\
\int e^x \, dx &=  e^x + C \\
\int b^x \, dx &= \dfrac{1}{\ln{b} \, b^x} + C  \\
\int  \dfrac{1}{\sqrt{1-x^2}}\, dx &= \arcsin{x} + C \\
\int  \dfrac{1}{1+x^2} \, dx &= \arctan{x} + C
\end{align*}

The integration sign $\int$ now has two distinct meanings.  
%
\mar{The \emph{definite} integral\\ is a number\ldots}
%
Originally, it was used to describe the {\em number\/}
	$$
	\int_a^b f(x) \, dx,
	$$
which was always calculated as the limit of a sequence of Riemann
sums.  Because this integral has a definite numerical value, it is
called the {\bf definite integral}\index{definite
  integral}\index{integral!definite}.  In its new meaning, the
integration sign is used to describe the antiderivative
	$$
	\int f(x) \, dx,
	$$
which is a {\em function}, not a number.  
%
\mar{\ldots while the \emph{indefinite} integral is a function}
%
To contrast the new use of $\int$ with the old, and to remind us that
the new expression is a variable quantity, it is called the {\bf
  indefinite integral}\index{indefinite integral}\index{integral!indefinite}.  The function that appears in either a definite or an
indefinite integral is called the {\bf integrand}.\index{integrand}
The terms ``antiderivative'' and ``indefinite integral'' are
completely synonymous. We will tend to use the former term in general
discussions, using the latter term when focusing on the process of
finding the antiderivative.

Because an indefinite integral represents an antiderivative, the
process of finding an antiderivative is sometimes called {\bf
  integration}\index{integration}.  Thus the term {\em integration\/},
as well as the symbol for it, has two distinct meanings.
	
	
\subsection{Using Antiderivatives}

According to the fundamental theorem, we can use an {\em indefinite\/}
integral to find the value of a {\em definite\/} integral---and this
largely explains the importance of
antiderivatives\index{antiderivative!importance of}.  In the
language of indefinite integrals, the statement of the fundamental
theorem in the box on page~\pageref{fund th-antider} takes the
following form.  \index{fundamental theorem of calculus}
\begin{center}
\begin{picture}(340,45)
\put(0,0){\framebox(340,45)
{${\displaystyle \int_a^b f(x) \, dx = F(b) - F(a)}$, where 
	$F(x) = {\displaystyle \int f(x) \, dx}$.
}}
\end{picture}
\end{center}

\medskip
{\bf Example 1}\quad  Find $\dint_1^4 x^2 \, dx$.  We have that
	$$
	\int x^2 \, dx = \tfrac{1}{3} x^3 + C;
	$$
it follows that
	$$
	\int_1^4 x^2 \, dx = \tfrac{1}{3} 4^3 + C - 
		\left(\tfrac{1}{3}1^3 +C)\right) 
		= \dfrac{64}{3} + C - \dfrac{1}{3} - C = 21.
	$$

\medskip
\noindent
{\bf Example 2}\quad Find $\dint_0^{\pi/2} \cos{t} \, dt$.  This time
the indefinite integral we need is
	$$
	\int \cos{t} \, dt = \sin{t} + C.
	$$
The value of the definite integral is therefore
	$$
	\int_0^{\pi/2} \cos{t} \, dt = \sin{\pi/2} + C - 
		\left( \sin{0} + C \right) = 1 + C - 0 - C = 1.
	$$
In each example the two appearances of $C$ cancel each other.  
%
\mar{The calculation\\ doesn't depend on $C$,\\ so take $C = 0$}
%
Thus $C$ does not appear in the final result.  This implies that it
does not matter which value of $C$ we choose to do the calculation.
Usually, we just take $C = 0$.

\medskip

\noindent
{\bf Notation}\index{notation!indefinite integral}\index{indefinite
  integral!notation}: Because expressions like $F(b) - F(a)$ occur
often when we are using indefinite integrals, we use the abbreviation
	$$
	F(b) - F(a) = 
		\left. F(x) \rule{0pt}{14pt}\,
		\right|_a^b.
	$$
Thus,
	$$
	\int_1^4 x^2 \, dx = \left. \dfrac{x^3}{3} \, \right|_1^4
		= \dfrac{64}{3} - \dfrac{1}{3} = 21.
	$$

\medskip

There are clear advantages to using antiderivatives to evaluate
definite integrals: we get exact values and we avoid many lengthy
calculations.  The difficulty is that the method works only if we can
find a formula for the antiderivative.

There are several reasons why we might not 
%
\mar{We may not recognize the \mbox{antiderivative\ldots}}	
%
find the formula we need.  For instance, the antiderivative we want
may be a function we have never seen before.  The function
$\arctan{x}$ is an example.

Even if we have a broad acquaintance with functions, we may still not
be able to find the formula for a given antiderivative.
%
\mar{\ldots and there may even be no formula for it}	
%
The reason is simple: for most functions we can write down, the
antiderivative is just not among the basic functions of calculus.  For
example, none of the basic functions, in any form or combination,
equals
	$$
	\int e^{x^2} dx \qquad \mbox{or} \qquad
		\int \dfrac{\sin{x}}{x} \, dx.
	$$
This does not mean that $e^{x^2}$, for example, has no
antiderivative.  On the contrary, the accumulation function
	$$
	\int_0^x e^{t^2} dt 
	$$
is an antiderivative of $e^{x^2}$.  It can be evaluated, graphed, and
analyzed like any other function.  What we lack is a {\em formula\/}
for this antiderivative in terms of the basic functions of calculus.


\subsection{Finding Antiderivatives}

In the rest of this chapter we will be deriving a number of statements
involving antiderivatives.  It is important to remember what such
statements mean.
\begin{center}
\begin{picture}(320,66)
\put(0,0){\framebox(320,66)
{
\begin{minipage}{280pt}
The following statements are completely equivalent: 
\linebreak
\vspace*{-6pt}
$$\int f(x)\,dx = F(x)+C \mbox{\qquad and \qquad} F' = f.$$
\end{minipage}
}}
\end{picture}
\end{center}
In other words, a statement 
%
\mar{We differentiate\\ to verify a statement\\ about antiderivatives} 
%
about antiderivatives can be verified by looking at a statement about
derivatives.  If it is claimed that the antiderivative of $f$ is $F$,
you check the statement by seeing if it is true that $F'=f$.  This was
how we verified the elementary antiderivatives we've considered so
far. Another way of expressing these relationships is
\[
\left(\int f(x)\,dx\right)' = f(x)\quad\mbox{and}\quad \int F'(x)\,dx
  = F(x)+C
\]
for any functions $f$ and $F$.


This duality between statements about derivatives and statements about
antiderivatives also holds when applied to more general statements.
Many of the most useful techniques for finding antiderivatives are
based on converting the general rules for taking derivatives of sums,
products, and chains into equivalent antiderivative form.
%
\mar{The constant multiple and addition rules}
%
We'll start with the simplest combinations, which involve a {\bf
  function multiplied by a constant} and a {\bf sum of two functions}.

\newpage

\mbox{}\vspace{-20pt}
\addtolength{\arraycolsep}{-3pt}
	$$
\begin{array}{rclcrcl}
	\multicolumn{3}{c}{\mbox{derivative form}} & &
		\multicolumn{3}{c}{\mbox{antiderivative form}}
		\\ [2pt]
	\hline \\ [-10pt]
	(k \cdot F)' &=& k \cdot F'& & 
		\dint k \cdot f \, dx &=& k \cdot \!\!\dint f \, dx \\[.15in]
	(F + G)' &=& F' + G' & \qquad\qquad & 
		\dint (f + g) \, dx &=& \dint f \, dx + 
		\dint g \, dx 
\end{array} 
	$$
\addtolength{\arraycolsep}{3pt}%
Let's verify these rules.  Set $\int f=F$ and $\int g = G$.  Then the
first rule is claiming that $\int (k\cdot f) = k\cdot F$---the
antiderivative of a constant times a function equals the constant
times the antiderivative of the function.  To verify this, we have to
show that when we take the derivative of the right-hand side, we get
the function under the integral on the left-hand side.  But $(k\cdot
F)'=k\cdot F'$ (by the derivative rule), which is just $k\cdot f $,
which is what we had to show.

Similarly, to show that $\int (f+g) = F + G$---the antiderivative of
the sum of two functions is the sum of their separate
antiderivatives---we differentiate the right-hand side and find
$(F+G)'=F'+G'$ (by the derivative rule for sums) which is just $f+g$,
so the rule is true.
\medskip

\medskip
\noindent
{\bf Example 3}\quad This example illustrates the use of both the
addition and the constant multiple rules.
\begin{align*}
\int (7e^x + \cos{x}) \, dx 
   &= \int 7e^x \, dx + \int \cos{x} \, dx \\
   &= 7 \int e^x \, dx + \int \cos{x} \, dx \\[6pt]
   &= 7 e^x + \sin{x} + C.
\end{align*}
To verify this answer, you should take the derivative of the
right-hand side to see that it equals the integrand on the left-hand
side.

\medskip
In the following sections we will develop the antidifferentiation
rules that correspond to the product rule and to the chain rule.  They
are called {\em integration by parts\/}\index{integration by
  parts}\index{integration by parts!correspondence to
  product rule} and {\em integration by
  substitution}\index{integration by substitution}\index{integration
  by substitution!correspondence to chain rule},
respectively.

While antiderivatives can be hard to find,
%
\mar{Trial and error}\index{trial and error}
%
they are easy to check.  This makes ``trial and error'' a good
strategy.  In other words, if you don't immediately see what the
antiderivative should be, but the function doesn't look too bad, try
guessing.  When you differentiate your guess, what you see may lead
you to a better guess.

\noindent
{\bf Example 4}\quad \label{`ex4'}  Find 
$$F(x) = \dint \cos(3x) \, dx.$$ 
Since the derivative of $\sin{u}$ is $\cos{u}$, it is reasonable to
try $\sin(3x)$ as an antiderivative for $\cos(3x)$.  Therefore:
\begin{align*}
\mbox{\sc first guess: } F(x) &= \sin(3x), \\
\mbox{\sc check: } F'(x) &= \cos(3x) \cdot 3 \neq \cos(3x).
\end{align*}
We wanted $F' = \cos(3x)$ but we got $F' = 3\cos(3x)$.  The chain rule
gave us an extra---and unwanted---factor of 3.  We can compensate for
that factor by multiplying our first guess by 1/3.  Then
\begin{align*}
\mbox{\sc second guess: } F(x) &= \tfrac{1}{3}\sin(3x), \\
\mbox{\sc check: } F'(x) &= \tfrac{1}{3}\cos(3x)\cdot 3  = 
		\cos(3x).
\end{align*}
Thus $\dint \cos(3x) \, dx = \tfrac{1}{3}\sin(3x)$.

\medskip

Because indefinite integrals are difficult to calculate, 
%
\mar{Tables of Integrals}\index{integrals!table of}
%
reference manuals in mathematics and science often include tables of
integrals.  There are sometimes many hundreds of individual formulas,
organized by the type of function being integrated.  A modest
selection of such formulas can be found at the back of this book.  You
should take some time to learn how these tables are arranged and get
some practice using them. You should also check some of the more
unlikely looking formulas by differentiating them to see that they
really are the antiderivatives they are claimed to be.

Computers are having a major impact on integration techniques.  We saw
in the last chapter that any continuous function---even one given by
the output from some laboratory recording device or as the result of a
numerical technique like Euler's method---can be approximated by a
polynomial (usually using lots of computation!), and the
antiderivative of a polynomial is easy to find.

Moreover, computer software packages 
%
\mar{Integration can\\ now be done\\ quickly and efficiently\\ by computer software}
%
which can find any existing formula for a definite integral are
becoming widespread and will probably have a profound impact on the
importance of integration techniques over the next several years.
Just as hand-held calculators have rendered obsolete many traditional
arts---like using logarithms for performing multiplications or knowing
how to interpolate in trig tables---there is likely to be a decreased
importance placed on humans' being adept at some of the more esoteric
integration techniques.  While some will continue to derive pleasure
from becoming proficient in these skills, for most users it will
generally be much faster, and more accurate, to use an appropriate
software package.  Nevertheless, for those going on in mathematics and
the physical sciences, it will still to be useful to be able to
perform some of the simpler integrations by hand reasonably rapidly.
The subsequent sections of this chapter develop the most commonly
needed techniques for doing this.


\subsection{Exercises}

\vspace{-10pt}

\num What is the inverse of the function $y = 1/x$?  Sketch the graph
of the function and its inverse.

\num What is the inverse of the function $y = 1/x^3$?  Sketch the
graph of the the function and its inverse.  Do the same for $y = 3x -
2$.

\numa Let $\theta=\arctan{x}$.  Then $\tan{\theta} = x$.  Refer to the
picture on page~\pageref{picture2} showing the relationship between $x$
and $\theta$.  Use this drawing to show that $\arctan(x) +
\arctan(1/x)$ is constant---that is, its value doesn't depend on $x$.
What is the value of the constant?  

\sub Use part (a), and the derivative of $\arctan{x}$, to find the
derivative of $\arctan(1/x)$.  

\sub Use the chain rule to verify your answer to part (b).

\num The logarithm for the base $b$ is defined as the inverse to the
exponential function with base $b$:
	$$
	y = \log_b{x} \qquad \mbox{if} \qquad x = b^y.
	$$
Using only the fact that $dx/dy = \ln{b} \cdot b^y$, deduce the formula
	$$
	\dfrac{dy}{dx} = \dfrac{1}{\ln{b}} \cdot \dfrac{1}{x}.
	$$
Note: this is purely an algebra problem; you don't need to invoke any
differentiation rules.

\numa Define $\arccos{x}$,\index{function!arccosine} the inverse of
the cosine function.  Be sure to limit the domain of the cosine
function to an interval on which it is one-to-one.

\sub Sketch the graph $y = \arccos{x}$.  How did you limit the range
of $y$?

\sub  Determine $dy/dx$. 

\numa If $\theta = \arcsin{x}$, refer to the picture on
page~\pageref{picture} reflecting the relation between $\theta$ and
$x$.  Using this picture, proceed as in problem 3 to to show that the
sum $\arcsin{x} + \arccos{x}$ is constant.  What is the value of the
constant?

\sub Use part (a), and the derivative of $\arcsin{x}$, to determine
the derivative of $\arccos{x}$.  Does this result agree with what you
got in the last exercise?

\num  Find a formula for $\dint \dfrac{dx}{\sqrt{1 - x^2}}$.


\num Verify that the antiderivatives given in the table on
page~\pageref{table} are correct.

\num Find an antiderivative of each of the following functions.  Don't
hesitate to use the ``trial and error'' method of Example 4 above.
\begin{center}
	\begin{tabular}{llll}
	$3$ 	& $5t$ 	& $-5t$ 	& $3 - 5t$ \\ [6pt]
	$7x^4$	& $\dfrac{1}{y^3}$	& 
		$e^{\displaystyle 2z}$ & 
		$u + \dfrac{1}{u}$ \\ [12pt]
	$(1 + w^3)^2$ \rule{15pt}{0pt} & 
		$\cos(5v)$ \rule{15pt}{0pt} & 
		$x^9 + 5x^7 - 2x^5$ \rule{15pt}{0pt} & 
		$\sin{t} \, \cos{t}$ 
	\end{tabular}
\end{center}

\num  Find a formula for each of the following indefinite integrals.

\medskip

\noindent
\begin{minipage}{150pt}

\sub $\dint 3x \, dx$

\sub $\dint 3u \, du$

\sub $\dint e^z \, dz$

\sub $\dint 5t^4 \, dt$

\sub $\dint 7y + \dfrac{1}{y} \, dy$

\sub $\dint 7y - \dfrac{4}{y^2} \, dy$
\end{minipage}
\hfill
\begin{minipage}{150pt}
		
\sub $\dint 5\sin{w} - 2\cos{w} \, dw$
	
\sub $\dint dx$

\sub $\dint e^{z+2} \, dz$

\sub $\dint \cos(4x) \, dx$

\sub $\dint \dfrac{5}{1+r^2}\,dr$

\sub $\dint \dfrac{1}{\sqrt{1-4s^2}}\, ds$
\end{minipage}
\hfill
\mbox{}


\numa Find an antiderivative $F(x)$ of $f(x) = 7$ for which $F(0) =
12$.

\sub Find an antiderivative $G(x)$ of $f(x) = 7$ for which $G(3) = 1$.

\sub Do $F(x)$ and $G(x)$ differ by a constant?  If so, what is the
value of that constant?


\numa Find an antiderivative $F(t)$ of $f(t) = t + \cos{t}$ for which
$F(0) = 3$.

\sub Find an antiderivative $G(t)$ of $f(t) = t + \cos{t}$ for which
$G(\pi/2) = -5$.

\sub Do $F(t)$ and $G(t)$ differ by a constant?  If so, what is the
value of that constant?

\num Find an antiderivative of the function $a + by$ when $a$ and $b$
are fixed constants.

\numa Verify that $(1 + x^3)^{10}$ is an antiderivative of $30x^2(1 +
x^3)^9$.

\sub  Find an antiderivative of $x^2(1 + x^3)^9$.

\sub  Find an antiderivative of $x^2 + x^2(1 + x^3)^9$.

\numa  Verify that $x \ln{x}$ is an antiderivative of $1 + \ln{x}$.

\sub Find an antiderivative of $\ln{x}$.  [Do you see how you can use
  part (a) to find this antiderivative?]


\num Recall that $F(y) = \ln(y)$ is an antiderivative of $1/y$ for
$y>0$.  According to the text, {\em every\/} antiderivative of $1/y$
over this domain must be of the form $\ln(y) + C$ for an appropriate
value of $C$.

\sub Verify that $G(y) = \ln(2y)$ is also an antiderivative of $1/y$.

\sub  Find $C$ so that $\ln(2y) = \ln(y) + C$.

\num Verify that $-\cos^2{t}$ is an antiderivative of $2\sin{t} \,
\cos{t}$.  Since you already know $\sin^2{t}$ is an antiderivative,
you should be able to show
	$$
	-\cos^2{t} = \sin^2{t} + C 
	$$
for an appropriate value of $C$.  What is $C$?

\num Since\label{log} the function $\ln{x}$ is defined only when $x >
0$, the equation
	$$
	\int \dfrac{1}{x} \, dx = \ln{x} + C
	$$
applies only when $x > 0$.  However, the integrand $1/x$ is defined
when $x < 0$ as well.  Therefore, it makes sense to ask what the
integral (i.e., antiderivative) of $1/x$ is when $x < 0$.

\sub When $x < 0$ then $ -x > 0$ so $\ln(-x)$ is defined.  In these
circumstances, show that $\ln(-x)$ is an antiderivative of $1/x$.

\sub Now put these two ``pieces'' of antiderivative together by
defining the function
\[
F(x) = \begin{cases}
\ln(-x)  &\mbox{if $x < 0$}, \\
\ln(x)   &\mbox{if $x > 0$}
\end{cases}
\]
Sketch together the graphs of the functions $F(x)$ and $1/x$ in such a
way that it is clear that $F(x)$ is an antiderivative of $1/x$.

\sub Explain why $F(x) = \ln{|x|}$.  For this reason a table of
integrals often contains the entry
	$$
	\int \dfrac{1}{x} \, dx = \ln{|x|} + C.
	$$

\sub Every function $\ln{|x|} + C$ is an antiderivative of $1/x$, but
there are even more.  As you will see, this can happen because the
domain of $1/x$ is broken into two parts.  Let
\[
G(x) = \begin{cases}
\ln(-x) & \mbox{if $x < 0$}, \\
\ln(x) + 1 & \mbox{if $x > 0$}.
\end{cases}
\]
Sketch together the graphs of the functions $G(x)$ and $1/x$ in such a
way that it is clear that $G(x)$ is an antiderivative of $1/x$.

\sub  Explain why there is no value of $C$ for which
	$$
	\ln{|x|} + C = G(x).
	$$
This shows that the functions $\ln{|x|} + C$ do not exhaust the set of
antiderivatives of $1/x$.

\sub Construct two more antiderivatives of $1/x$ and sketch their
graphs.  What is the general form of the new antiderivatives you have
constructed?  (A suggestion: you should be able to use two separate
constants $C_1$ and $C_2$ to describe the general form.)

\num On page~\pageref{tangent} of the text there is an antiderivative
for the tangent function:
	$$
	\int \tan{x} \, dx = - \ln(\cos{x}).
	$$
However, this is not defined when $x$ makes $\cos{x}$ either zero or
negative.

\sub How many separate intervals does the domain of $\tan{x}$ break
down into?

\sub For what values of $x$ is $\cos{x}$ equal to zero, and for what
values is it negative?

\sub Modify the antiderivative $-\ln(\cos{x})$ so that it {\em is\/}
defined when $\cos{x}$ is negative.  (How is this problem with the
logarithm function treated in the previous question?)

\sub  In a typical table of integrals you will find the statement
	$$
	\int \tan{x} \, dx =  - \ln{|\cos{x}|} + C.
	$$
Explain why this does not cover all the possibilities.

\sub Give a more precise expression for $\dint \tan{x}\,dx$, modelled
on the way you answered part (f) of problem 18.  How many different
constants will you need?

\sub Find a function $G$ that is an antiderivative for $\tan{x}$ and
that also satisfies the following conditions:
\[
G(0) = 5,  \qquad G(\pi) = -23, \qquad G(17\pi) = 197.
\]

\num In the table on page~\pageref{table} the antiderivative of $x^p$
is given as 
\[
\dfrac{1}{p+1} x^{p+1} +C.
\]
For some values of $p$ this is correct, with only a single constant
$C$ needed.  For other values of $p$, though, the domain of $x^p$ will
consist of more than one piece, and $\dfrac{1}{p+1} x^{p+1} $ can be
modified by a different constant over each piece.  For what values of
$p$ does this happen?

\num  Find $F'(x)$ for the following functions.  In parts (a), (b), and (d) do the problems two ways: by finding an antiderivative, and by using the fundamental theorem to get the answer without evaluating an antiderivative.  Check that the answers agree.

\sub  $F(x) = \dint_0^x (t^2 + t^3) \, dt$. 

\sub  $F(x) = \dint_1^x \dfrac{1}{u} \, du$.

\sub  $F(x) = 
	\dint_1^x \dfrac{v}{1+v^3} \, dv$.  

\sub  $F(x) = \dint_0^{x^2} \cos{t} \, dt$.  

\sub  $F(x) = 
	\dint_1^{x^2} \dfrac{v}{1+v^3} \, dv$. 
	\quad [Hint: let $u = x^2$ and use the chain rule.]
\smallskip

\noindent
Comment: It may seem that parts (c) and (e) are more difficult than
the others.  However, there is a way to apply the fundamental theorem
of calculus here to get answers to parts (c) and (e) quickly and with
little effort.

\num  Consider the two functions
	$$
	F(x) = \sqrt{1 + x^2} - 1 \quad \mbox{and} \quad
		G(x) = \int_0^x \dfrac{t}{\sqrt{1+t^2}} \, dt.
	$$

\sub  Show that $F$ and $G$ both satisfy the initial value problem
	$$
	y' = \dfrac{x}{\sqrt{1+x^2}},  \qquad y(0) = 0.
	$$

\sub Since an initial value problem typically has a {\em unique\/}
solution, $F$ and $G$ should be equal.  Assuming this, determine the
exact value of the following definite integrals.
	$$
	\int_0^1 \dfrac{t}{\sqrt{1+t^2}} \, dt, \qquad
		\int_0^2 \dfrac{t}{\sqrt{1+t^2}} \, dt, \qquad 
		\int_0^5 \dfrac{t}{\sqrt{1+t^2}} \, dt. 
	$$

\num The connection between integration and differentiation that is
provided by the fundamental theorem of calculus makes it possible to
determine an integral by solving a differential equation.  For
example, the accumulation function
	$$
	A(x) = \int_0^x e^{-t^2} dt 
	$$
is the solution to the initial value problem
	$$
	y' = e^{-x^2} \, , \qquad y(0) = 0.
	$$
Therefore, $A(x)$ can be found by solving the differential equation.
As you have seen, Euler's method is a useful way to solve differential
equations.

\sub Use either a program (e.g., PLOT) or a differential equation
solver on a computer to get a graphical solution $A(x)$ to the initial
value problem above.

\sub Sketch the graph of $y = A(x)$ over the domain $0 \leq x \leq 4$.

\sub Your graph should increase from left to right.  How can you tell
this even before you see the computer output?

\sub Your graph should level off as $x$ increases.  Determine $A(5)$,
$A(10)$, $A(30)$.  (Approximations provided by the computer are
adequate here.)

\sub Estimate ${\displaystyle \lim_{x \to \infty} A(x)}$.  \hfill [The
  {\em exact\/} value is $\sqrt{\pi}/2$.]

\sub Determine $\dint_0^1 e^{-t^2} dt$ and $\dint_0^2 e^{-t^2} dt$.

\sub  Determine $\dint_1^2 e^{-t^2} dt$.

\num Find the area under the curve $y = x^3 + x$ for $x$ between 1
and~4.  (See chapter 6.3.)

\num Find the area under the curve $y = e^{3x}$ for $x$ between 0 and
$\ln{3}$.

\num The {\bf average value} of the function $f(x)$ on the interval $a
\leq x \leq b$ is the integral
\[
\dfrac{1}{b-a} \int_a^b f(x) \, dx.
\]
(See the discussion of average value in chapter~6, pages
\pageref{average-begin}--\pageref{average-end}.)

\sub Find the average value of each of the functions $y = x$, $x^2$,
$x^3$, and $x^4$ on the interval $0 \leq x \leq 1$.

\sub Explain, using the graphs of $y = x$ and $y = x^2$, why the
average value of $x^2$ is less than the average value of $x$ on the
interval $[0, 1]$.

\numa  What are the maximum, minimum, and average values of the function 
	$$
	f(x) = x + 2 e^{-x}
	$$
on the interval $[0, 3]$?

\sub Sketch the graph of $y = f(x)$ on the interval $[0, 3]$.  Draw
the line $y = \mu$, where $\mu$ is the average value you found in part
(a).

\sub For which $x$ does the graph of $y = f(x)$ lie above the line $y
= \mu$, and for which $x$ does it lie below the line?  The region
between the graph and the line has two parts---one is above the line
(and below the graph) and the other is below the line (and above the
graph).  Shade these two regions and compare their areas: which is
larger?  \hfill [The two are equal.]


%%%%%%%%%%%%%%%%%%%%%%%%%%%%%%%%%%%%%%%%%%%%%%%%%%%%%%%%%%%%%%%%%%%%%%%%
%                                                                      %
%                               Section 2                              %
%                                                                      %
%%%%%%%%%%%%%%%%%%%%%%%%%%%%%%%%%%%%%%%%%%%%%%%%%%%%%%%%%%%%%%%%%%%%%%%%


\section{Integration by Substitution}

In the preceding section we converted a couple of general rules for
differentiation---the rule for the derivative of a constant times a
function and the rule for the derivative of the sum of two
functions---into equivalent rules in integral form.  In this section
we will develop the integral form of the chain rule and see some of
the ways this can be used to find antiderivatives.

Suppose we have functions $F$ and $G$, with corresponding derivatives
$f$ and $g$.  Then the chain rule says
$$ 
\left(F(G(x))\right)'=F'(G(x))\,G'(x)=f(G(x))\,g(x).
$$
If we now take the indefinite integral of these equations, we get
$$
F(G(x))+C=\int (F(G(x)))'\,dx=\int f(G(x))\,g(x)\,dx ,
$$
where $C$ can be any constant.

Reversing these equalities 
%
\mar{The integral form\\ of the chain rule}
%
to get a statement about integrals, we obtain:\label{subrule}
\begin{center}
\begin{picture}(240,40)
	\begin{thicklines}
	\put(0,0){\framebox(240,40)
		{\bf \boldmath $\displaystyle{\int f(G(x))\,g(x)\,dx=F(G(x)) + C.}$}}
	\end{thicklines}
\end{picture}
\end{center}

This somewhat unpromising expression turns out to be surprisingly
useful.  Here's how: Suppose we want to find an indefinite integral,
and see in the integrand a {\em pair} of functions $G$ and $g$, where
$G'=g$ and where $g(x)$ can be factored out of the integrand.  
%
\mar{Reduction methods transform problems into equivalent, simpler, problems} 
%
We then find a function $f$ so that the integrand can be written in
the form $f(G(x))g(x)$.  Now we only have to find an antiderivative
for $f$.  Once we have such an antiderivative, call it $F$, then the
solution to our original problem will be $F(G(x))$.  Thus, while our
original antiderivative problem is not yet solved, it has been reduced
to a different, simpler antiderivative problem that, when all goes
well, will be easier to evaluate.  Such {\bf reduction methods} are
typical of many integration techniques.  We will see other examples in
the remainder of this chapter.

\medskip
{\bf Example 1}\quad  Suppose we try to find a formula for the integral
	$$
	\int 3x^2 \left(1 + x^3\right)^7 dx.
	$$
One way would be to multiply out the expression $\left(1 +
x^3\right)^7$, making the integrand a polynomial with many separate
terms of different degrees.  (The highest degree would be 23; do you
see why?)  We could then carry out the integration ``term-by-term,''
using the rules for sums and constant multiples of powers of $x$ that
were given in the previous section.  But this is tedious---even
excruciating.

Instead, notice that the expression $3x^2$ is the derivative of
$1+x^3$.  If we let $G(x)=1+x^3$ and $g(x)=3x^2$, we can then write
the integrand as
	$$
	3x^2 \left(1 + x^3\right)^7=(G(x))^7g(x).
	$$
Now $(G(x))^7$ is clearly just $f(G(x))$ where $f$ is the function
which raises its input to the 7th power---$f(x)=x^7$ for any input
$x$.  But we recognize $f$ as an elementary function whose
antiderivative we can write down immediately as
$F(x)=\tfrac{1}{8}x^8$.  Thus the solution to our original problem
will be
\[
\int 3x^2 \left(1 + x^3\right)^7 dx = F(G(x))+C 
         = \dfrac{1}{8}\left(1 + x^3\right)^8+C.
\]
As with any integration problem, we can check our answer by taking the
derivative of the right-hand side to see if it agrees with the
integrand on the left.  You should do this whenever you aren't quite
sure of your technique (or your arithmetic!).

Note that the term $3x^2$ that appeared in the integrand above was essential for the procedure to work.  The integral
	$$
	\int \left(1 + x^3 \right)^7 dx
	$$
{\em cannot\/} be found by substitution, even though it appears to
have a simpler form.  (Of course, the integral {\em can\/} be found by
multiplying out the integrand.)

%\medskip
\newpage

\noindent
{\bf Using differential notation}\index{differential
  notation}\label{differentials}\quad 
%
\mar{A compact notation for expressing the integral form of the chain rule} 
%
So far the symbol $dx$ (the {\bf differential} of~$x$) under the
integral sign has simply been an appendage, tagging along to suggest
the $\Delta x$ portion of the Riemann sums approximating definite
integrals.  It turns out we can take advantage of this notation to use
the integral form of the chain rule more compactly.

Instead of naming the functions $G$ and $g$ as above, we introduce a
new variable $u=G(x)$.  Then
$$
\dfrac{du}{dx} = G'(x)=g(x),
$$
and it is suggestive to multiply out this ``quotient'' to get
\[
du=g(x)\,dx.
\]  
While this is reminiscent of the microscope equation we met in
chapter~3, and 18th century mathematicians took this equation
seriously as a relation between two ``infinitesimally small''
quantities $dx$ and $du$, we will view it only as a convenient
mnemonic device. To see how this simplifies computations, reconsider
the previous example.  If we let $u=1+x^3$, then $du=3x^2\,dx$, so we
can write
	$$
	3x^2 (1 + x^3)^7 \, dx = 
		(\underbrace{1 + \rule[-3pt]{0pt}{6pt}x^3}_
			{\displaystyle u})^7 
		\underbrace{3x^2 \rule[-3pt]{0pt}{6pt}\, dx}_
			{\displaystyle du} =
		u^7 \, du.
	$$
It follows that
\begin{align*}
\int 3x^2 \left(1 + x^3 \right)^7 dx &= \int u^7 \, du \\
  &= \tfrac{1}{8} u^8 + C \\
  &= \tfrac{1}{8} \left(1 + x^3\right)^8 + C,
\end{align*}
as before.  We have arrived at the same answer without having to
introduce the cumbersome language of all the auxiliary functions---we
simply {\bf substituted} the variable $u$ for a certain expression in
$x$ (which we called $G(x)$ before), and replaced $G'(x)dx$ by
$du$. For this reason this technique is called {\bf integration by
  substitution}.  You should always be clear, though, that integration
by substitution is just the integral form of the chain rule, a
relationship that becomes clear whenever you check the answer
substitution gives you.


\medskip
\noindent
{\bf Example 2}\quad  Can we use the method of substitution to find
	$$
	\int \dfrac{e^{5x}}{6 + e^{5x}} \, dx?
	$$
The numerator is almost the derivative of the denominator.  This
suggests we let $G(x) = 6+e^{5x}$, giving $g(x)=G'(x)= 5e^{5x}$.
Since we need to be able to factor $g(x)$ out of the integrand, we
multiply numerator and denominator by 5 to get
\begin{align*}
\int \dfrac{e^{5x}}{6 + e^{5x}} \, dx 
  &= \int \dfrac{1}{5} \cdot \dfrac{1}{6+e^{5x}}5e^{5x}\,dx \\
  &= \int \dfrac{1}{5} \cdot \dfrac{1}{G(x)}g(x)\,dx \\
  &= \dfrac{1}{5}\int f(G(x))g(x)\, dx,
\end{align*}
where $f$ is just the reciprocal function---$f(x)=1/x$.  But an
antiderivative for $f$ is just $\ln{x}$, so the desired antiderivative
is just
$$
\dfrac{1}{5}F(G(x))+C= \dfrac{1}{5}\ln(6+e^{5x})+C.
$$
As usual, you should check this answer by differentiating to see that
you really do get the original function.

Now let's see how this works using differential notation.  If we set
$u = 6 + e^{5x}$, then
	$$
	\dfrac{du}{dx} = 5e^{5x} \qquad \mbox{and} \qquad
		du = 5e^{5x} \, dx.
	$$
Again we insert a factor of 5 in the numerator and an identical one in
the denominator to balance it.  Substitutions for $u$ and $du$ then
yield the following:
\begin{align*}
\int \dfrac{1}{5}\cdot \dfrac{5e^{5x}}{6 + e^{5x}} \, dx 
  &= \dfrac{1}{5} \int \dfrac{5e^{5x} \, dx}{6 + e^{5x}} \\[3pt]
  &= \dfrac{1}{5} \int \dfrac{du}{u} \\[4pt]
  &= \dfrac{1}{5} \ln(u) + C \\[3pt]
  &= \dfrac{1}{5} \ln(6+e^{5x}) + C, 
\end{align*}
\smallskip
as before.

\newpage
%\medskip
	
The two examples above have the same structure.  
%
\mar{The basic structure}	
%
In both, a certain function of $x$ is selected and called $u$; part of
the integrand, namely $u'\,dx$, becomes $du$, the rest becomes one of
the basic functions of $u$.  Specifically:
	$$
	\begin{array}{ccccccc}
	\mbox{integrand} && u && du 
		&& \mbox{function of $u$} \\[2pt] \hline
	3x^2 \left(1 + x^3\right)^7 dx && 1 + x^3 
		&& 3x^2 \rule{0pt}{18pt}\, dx	&& u^7 \\[4pt]
	\dfrac{e^{5x}}{6 + e^{5x}} \, dx 
		&& 6 + e^{5x} && 5e^{5x} \, dx 
		&& \dfrac{1}{5} \cdot \dfrac{1}{u}
	\end{array}
	$$



Note that you may have to do a bit of algebraic reshaping of the
integrand\index{integrand} to cast it in the proper form.  For
example, we had to insert a factor of 5 to the numerator of the second
example to make the numerator be the derivative of the denominator.
There is no set routine to be followed to find an antiderivative most
efficiently, or even any way to know whether a particular method will
work until you try it.  Success comes with experience and a certain
amount of intelligent fiddling until something works out.
\medskip

\noindent
{\bf Example 3}\quad The method of substitution is useful in simple
problems, too.  Consider
	$$
	\int \cos(3t) \, dt.
	$$
If we set $u = 3t$, then $du = 3 \, dt$ and
\begin{align*}
\int \cos(3t) \, dt 
  &= \int \cos(u) \cdot \tfrac{1}{3} \, du \\[2pt]
  &= \tfrac{1}{3} \int \cos(u) \, du \\[4pt]
  &= \tfrac{1}{3} \sin(u) + C \\[6pt]
  &= \tfrac{1}{3} \sin(3t) + C. 
\end{align*}

\subsection{Substitution in Definite Integrals}

Until now we have been using the technique of substitution to find
antiderivatives---that is, to evaluate {\em indefinite integrals}.
Refer back to the integral form of the chain rule given in the box on
page~\pageref{subrule}, and see what happens when we use this equation
to evaluate a {\em definite integral}.  Suppose we want to evaluate
$$
\int_a^bf(G(x))\,g(x)\,dx.
$$ 
We know that $F(G(x))$ is an antiderivative, so by the fundamental
theorem we have
$$
\int_a^bf(G(x))\,g(x)\,dx=\left. F(G(x))\right|_a^b=F(G(b))-F(G(a)).
$$

Now suppose we make the substitution $u=G(x)$ and $du=g(x)\,dx$.  Then
as $x$ goes from $a$ to $b$, $u$ will go from $G(a)$ to $G(b)$.
Moreover, we have
$$\int_{G(a)}^{G(b)} f(u)\,du=\left. F(u)\right|_{G(a)}^{G(b)} = F(G(b))-F(G(a)),$$
so the two definite integrals have the same value. In other words,
\label{defint}
\begin{center}
\framebox[380pt]
{\parbox{350pt}{\vspace{5pt}
\hfill \textbf{If we make the substitution \boldmath $u=G(x)$, then 
       \hfill \mbox{} \\[8pt]
$\dint_a^b f(G(x))\,g(x)\,dx = F(G(b))-F(G(a))=\dint_{G(a)}^{G(b)} f(u)\,du$.
}\vspace{5pt}
}}
\end{center}

This means that to evaluate a definite integral by substitution, we
can do everything in terms of $u$.  We don't ever need to find an
antiderivative for the original integrand in terms of $x$ or use the
original limits of integration.

\medskip
\noindent
{\bf Example 4}\quad Consider the definite integral
$$
\int_0^{\pi/2} \dfrac{\cos{x}\,dx}{1+\sin{x}}.
$$
We can evaluate this integral by making the substitution
$u=1+\sin{x}$, and $du=\cos{x}\,dx$.  Moreover, as $x$ goes from 0 to
$\pi/2$, $u$ goes from 1 to 2.  Therefore
$$
\int_0^{\pi/2} \dfrac{\cos{x}\,dx}{1+\sin{x}} 
   = \int_1^2 \dfrac{du}{u}=\left.\ln{u}\right|_1^2=\ln{2}.
$$
Check that this is the same answer you would have gotten if you had
expressed the antiderivative $\ln{u}$ in terms of $x$ and evaluated
the result at the limits on the original integral.

\medskip

\noindent
{\bf Example 5}\quad Evaluate
\[
\int_0^1 6 x^2(1+x^3)^4\,dx.
\]
With the substitution $u=1+x^3$ and $du=3x^2\,dx$, as $x$ goes from 0
to 1, $u$~goes from 1 to 2. Our integral thus becomes
\[
\int_0^1 6 x^2(1+x^3)^4\,dx=\int_0^2 2 u^4\,du
   =\left. \dfrac{2}{5}u^5 \, \right|_1^2 =\dfrac{64}{5}.
\]

\subsection{Exercises}

\vspace{-10pt}

\num Evaluate the following using substitution.  Do parts (a) through
(e) in two ways: {\rm i.}  by writing the integrand in the form
$f(G(x))g(x)$ (or $y$ or $t$ or whatever the variable is) for
appropriate functions $f$, $G$, and $g$, with $G'=g$, and then finding
$F = \int f$; and {\rm ii.} using differential notation.  Do the
remaining parts in the way you feel most confident.

\noindent
\begin{minipage}{150pt}
\sub ${\dint 2y(y^2+1)^{50}\,dy}$

\sub ${\dint \sin(5z)\,dz}$

\sub ${\dint \dfrac{e^{\sqrt x}}{\sqrt x}\,dx}$

\sub ${\dint (5t+7)^{50}\,dt}$

\sub ${\dint 3u^2 \, \root 3 \of {u^3+8}\,du}$

\sub ${\dint \dfrac{1}{2v+1}\,dv}$

\sub ${\dint \tan{x}\,dx}$

\sub ${\dint \tan^2(x)\sec^2(x)\,dx}$

\sub ${\dint \sec(x/2)\tan(x/2)\,dx}$
\end{minipage}
\hfill
\begin{minipage}{150pt}
\sub ${\dint \sin(w)\,\sqrt{\cos(w)}\,dw}$

\sub ${\dint \dfrac{\sin(\sqrt{s})}{\sqrt{s}}\,ds}$

\sub ${\dint \sqrt{3-x}\,dx}$

\sub ${\dint \dfrac{dr}{r\ln{r}}}$

\sub ${\dint e^x\sin(1+e^x)\,dx}$

\sub ${\dint \dfrac{y}{1+y^2}\,dy}$

\sub ${\dint \dfrac{w}{\sqrt{1-w^2}}\, dw}$

\sub ${\dint \dfrac{1}{1+4y^2}\,dy}$

\sub ${\dint \dfrac{1}{\sqrt{1-9w^2}}\, dw}$
\end{minipage}
\hfill
\mbox{}

\num Use integration by substitution to find the numerical value of
the following.  In four of these you should get your answer in two
ways: i) by finding an antiderivative for the given integrand, and ii)
by using the observation in the box on page~\pageref{defint},
comparing the results.  You should also check your results for three
of the problems by finding numerical estimates for the integrals using
RIEMANN\index{program!RIEMANN}.

\noindent
\begin{minipage}{150pt}
\sub ${\dint_0^1 \dfrac{e^s}{e^s+1}\,ds}$

\sub ${\displaystyle \int_0^{\ln{e}}
\dfrac{e^s}{e^s+1}\,ds}$

\sub ${\displaystyle \int_1^3 \dfrac{1}{2x+1}\,dx}$

\sub ${\displaystyle \int_{-3}^{-1} \dfrac{1}{2x+1}\,dx}$
\end{minipage}
\hfill
\begin{minipage}{150pt}
\sub ${\displaystyle \int_0^1
	\dfrac{t}{\sqrt{1 + t^2}}\,dt}$

\sub ${\displaystyle \int_0^1
	\dfrac{\sin(\pi\sqrt{t})}{\sqrt{t}}\,dt}$

\sub ${\displaystyle \int_0^2
	\dfrac{1}{1+ (x^2/4)}\,dx}$

\sub ${\displaystyle \int_0^{\tfrac{1}{3}}
	\dfrac{1}{\sqrt{1-9y^2}}\,dy}$
\end{minipage}
\hfill
\mbox{}

\num  This question concerns the integral 
$I = \dint \sin{x} \, \cos{x} \, dx$.

\sub  Find $I$ by using the substitution $u = \sin{x}$.

\sub  Find $I$ by using the substitution $u = \cos{x}$.

\sub Compare your answers to (a) and (b).  Are they the same?  If not,
how do they differ?  Since both answers are antiderivatives of
$\sin{x} \, \cos{x}$, they should differ only by a constant.  Is that
true here? If so, what is the constant?

\sub  Now calculate the value of the {\em definite\/} integral
	$$
	\int_0^{\pi/2} \sin{x}\,\cos{x} \,dx
	$$
twice, using the two {\em indefinite\/} integrals you found in (a) and
(b).  Do the two values agree, or disagree?  Is your result consistent
with what you expect?


\numa  Find all functions $y = F(x)$ that satisfy the differential equation 
	$$
	\dfrac{dy}{dx} = x^2 \left( 1 + x^3 \right)^{13}.
	$$

\sub From among the functions $F(x)$ you found in part (a), select the
one that satisfies $F(0) = 4$.

\sub From among the functions $F(x)$ you found in part (a), select the
one that satisfies $F(-1) = 4$.


\num  Find a function $y = G(t)$ that solves the initial value problem
	$$
	\dfrac{dy}{dt} = t e^{-t^2}
		\qquad y(0) = 3.
	$$
	
\numa What is the average value of the function $f(x) =
x/\sqrt{1+x^2}$ on the interval $[0, 2]$?

\sub Show that the average value of the $f(x)$ on the interval $[-2,
  2]$ is 0.  Sketch a graph of $y = f(x)$ on this interval, and
explain how the graph also shows that the average is 0.

\numa Sketch the graph of the function $y = x e^{-x^2}$
on the interval $[0, 5]$.

\sub Find the area between the graph of $y = x e^{-x^2}$
and the $x$-axis for $0 \leq x \leq 5$.

\sub Find the area between the graph of $y = x e^{-x^2}$
and the $x$-axis for $0 \leq x \leq b$.  Express your answer in terms
of the quantity $b$, and denote it $A(b)$.  Is $A(5)$ the same number
you found in part(b)?  What are the values of $A(10)$, $A(100)$,
$A(1000)$?

\sub It is possible to argue that the area between the graph of $y = x
e^{-x^2}$ and the {\em entire\/} positive $x$-axis is
1/2.  Can you develop such an argument?


\numa  Use a computer graphing utility to establish that
	$$
	\sin^2{x} = \dfrac{1 - \cos(2x)}{2}.
	$$
Sketch these graphs.


\sub Find a formula for $\dint \sin^2{x} \, dx$.  (Suggestion: replace
$\sin^2{x}$ by the expression involving $\cos(2x)$, above, and
integrate by substitution.)

\sub What is the average value of $\sin^2{x}$ on the interval $[0,
  \pi]$?  What is its average value on any interval of the form $[0, k
  \pi]$, where $k$ is a whole number?

\sub Explain your results in part (c) in terms of the graph of
$\sin^2{x}$ you drew in part (a).

\sub Here's a differential equations proof of the identity in part
(a).  Let $f(x)=\sin^2{x}$, and let $g(x)=(1 - \cos(2x))/2$.  Show
that both of these functions satisfy the initial value problem
$$y''=2-4y \quad\mbox{with}\quad y(0)=0\quad\mbox{and}\quad y'(0)=0.$$
Hence conclude the two functions must be the same.  

\newpage
%%%%%%%%%%%%%%%%%%%%%%%%%%%%%%%%%%%%%%%%%%%%%%%%%%%%%%%%%%%%%%%%%%%%%%%%
%                                                                      %
%                               Section 3                              %
%                                                                      %
%%%%%%%%%%%%%%%%%%%%%%%%%%%%%%%%%%%%%%%%%%%%%%%%%%%%%%%%%%%%%%%%%%%%%%%%

\section{Integration by Parts}\index{integration by parts}

As in the previous section, 
%
\mar{An integral form\\ of the product rule}
%
suppose we have functions $F$ and $G$,
with corresponding derivatives $f$ and $g$.  If we use the product
rule\index{product rule} to differentiate $F(x) \cdot G(x)$, we get:
	$$
	\left( F \cdot G \right)' = 
		F \cdot G' + F' \cdot G = F\cdot g + f\cdot G.
	$$
We can turn this into a statement about indefinite
integrals\index{integral!indefinite}\index{indefinite integral}:
	$$
	\int \left( F \cdot g + f \cdot G \right) \, dx = 
		\int \left( F \cdot G \right)' \, dx = 
		F(x) \cdot G(x) + C.
	$$
Unfortunately, in this form the statement is not especially useful; 
it applies only when the integrand has two terms of the special form $f \cdot 	g' + f' \cdot g$.  However, if we rewrite the statement in the form
\begin{center}
\begin{picture}(230,40)
	\begin{thicklines}
	\put(0,0){\framebox(230,40)
		{\bf \boldmath $\dint F \cdot g \, dx = 
			F \cdot G - \dint f \cdot G \, dx$}}
	\end{thicklines}
\end{picture}
\end{center}
it becomes very useful.

\medskip

\noindent
{\bf Example 1}\quad  We will use the formula in the box to find
	$$
	\int x \cdot \cos{x} \, dx.
	$$
If we label the parts of this integrand as follows:
	$$
	F(x) = x \qquad \qquad g(x) = \cos{x},
	$$
then we have
	$$
	f(x) = 1 \qquad \mbox{and} \qquad G(x) = \sin{x}.
	$$
According to the formula,
\begin{align*}
\int x \cdot \cos{x} \, dx 
  &= x \cdot \sin{x} - \int \sin{x} \, dx \\
  &= x \cdot \sin{x} + \cos{x} + C.
\end{align*}

The integrand is first broken into two parts---in this case, 
%
\mar{Integrate only part\\ of the integrand}
%
$x$ and $\cos{x}$.  One part is differentiated while the other part is
integrated.  (The part we integrated is $g(x) = \cos{x}$, and we got
$G(x) = \sin{x}$.)  For this reason, the rule described in the box is
called {\bf integration by parts}\index{integration by parts}.

\newpage

%\medskip

As with integration by substitution, integration by parts exchanges
one integration task for another: Instead of finding an antiderivative
for $F \cdot g$, we must find one for $f \cdot G$.  The idea is to
``trade-in'' one integration problem for a more readily solvable one.

\medskip

\noindent
{\bf Example 2}\quad  Use integration by parts to find\label{lnder}
	$$
	\int \ln(x) \, dx.
	$$
At first glance we can't integrate by parts, because there aren't two parts!  But note that we can write
	$$
	\ln(x) = \ln(x) \cdot 1,
	$$
and then set 
	$$
	F(x) = \ln(x), \qquad \qquad g(x) = 1.
	$$
This implies 
	$$
	f(x) = \dfrac{1}{x} \qquad \mbox{and} \qquad
		G(x) = x,
	$$
and the integration by parts formula now gives us
\begin{align*}
\int \ln(x) \, dx 
  &= x \cdot \ln(x) - \int x \cdot \dfrac{1}{x} \, dx \\[4pt]
  &= x \cdot \ln(x) - \int 1 \, dx \\[4pt]
  &= x \cdot \ln(x) - x + C.
\end{align*}


We thus see that integration by parts\index{integration by
  parts}---like integration by substitution---is an art rather than a
set routine.
%
\mar{The ingredients\\ of a successful\\ integration by parts}	
%
If integration by parts is to work, several things must happen.
First, you need to see that the method might actually apply.  (In
Example 2 this wasn't obvious.)  Next, you need to identify the parts
of the integrand that will be differentiated and integrated,
respectively.  The wrong choices can lead you away from a solution,
rather than towards one.  (See example 3 below for a cautionary
tale\index{integration by parts!cautionary tale}.)  Finally,
you need to be able to carry out the integration of the new integral
$\int f \cdot G \, dx$.  As you work you may have to reshape the
integrand algebraically.  Technique comes with practice, and luck is
useful, too.
\medskip

\noindent
{\bf Example 3}\quad  Use integration by parts to find
	$$
	\int t \cdot e^t \, dt.
	$$
Set
	$$
	F(t) = e^t, \qquad \mbox \qquad g(t) = t;
	$$
then 
	$$
	f(t) = e^t, \qquad \mbox \qquad G(t) = \dfrac{t^2}{2}.
	$$
The integration by parts formula then gives
	$$
	\int t \cdot e^t \, dt =
		\dfrac{t^2}{2} e^t - \int \dfrac{t^2}{2} e^t\,dt.
	$$
While this is a true statement, we are not better off---the new
integral is {\em not\/} simpler than the original.  \mar{What went
  wrong?}  A solution is eluding us here.  You will have a chance to
do this problem properly in the exercises.


\subsection{Exercises}

\vspace{-10pt}

\num Use integration by parts to find a formula for each of the
following integrals.

\noindent
\begin{minipage}{150pt}
\sub ${\displaystyle \int x \sin{x}\,dx}$

\sub ${\displaystyle \int te^{t}\,dt}$

\sub ${\displaystyle \int we^{-w}\,dw}$

\sub ${\displaystyle \int x\ln{x}\,dx}$

\sub ${\displaystyle \int \arcsin{x}\,dx}$
\end{minipage}
\hfill
\begin{minipage}{150pt}
\sub ${\displaystyle \int \arctan{x}\,dx}$

\sub ${\displaystyle \int x^2e^{-x}\,dx}$ 

\sub $\dint u^2 \cos{u} \, du$    

\sub ${\displaystyle \int x\sec^2{x}\,dx}$

\sub ${\displaystyle \int e^{2x}(x + e^x)\,dx}$
\end{minipage}
\hfill
\mbox{}

\medskip

\noindent
(Suggestion for part (g): Apply integration by parts twice.  After the
first application you should have an integral that can itself be
evaluated using integration by parts.)


\num  Use integration by parts to obtain a formula for
	$$
	\int (\ln{x})^2\,dx.
	$$
Choose $f = \ln{x}$ and also $g' = \ln{x}$.  To continue you need to
find $g$, the antiderivative of $\ln{x}$, but this has already been
obtained in the text.


\numa  Find $\dint x^2 e^x \, dx$.

\sub  Find $\dint x^3 e^x \, dx$.  (Reduce this to part (a)).

\sub  Find $\dint x^4 e^x \, dx$.  

\sub  What is the general pattern here?  Find a formula for $\dint x^n e^x \, dx$, where $n$ is any positive integer.

\sub  Find $\dint e^x \left(5x^2 - 3x + 7\right) \, dx$.

\numa Draw the graph of $y=\arctan{x}$ over the interval $0\leq x\leq
1$.  You could have gotten the same graph by thinking of $x$ as a
function of $y$---write down this relationship and the corresponding
$y$ interval.

\sub Evaluate
$$\int_0^1 \arctan{x}\,dx$$
and show on your graph the area this corresponds to.

\sub Evaluate
$$\int_0^{\pi/4} \tan{y}\,dy$$
and show on your graph the area this corresponds to.

\sub If we add the results of part (b) and part (c), what do you get?
From the geometry of the picture, what should you have gotten?

\num Repeat the analysis of the preceding problem by calculating the
value of
\[
\int_0^2 x^3\,dx+\int_0^8 y^{1/3}\,dy,
\]
and seeing if it agrees with what you would predict by looking at the
graphs.

\newpage

\num Generalize the preceding two problems to the case where $f$ and
$g$ are any two functions that are inverses of each other whose graphs
pass through the origin.


\numa What is the average value of the function $\ln{x}$ on the
interval $[1, e]$?

\sub What is the average value of $\ln{x}$ on $[1, b]$?  Express this
in terms of $b$.  Discuss the following claim: The average value of
$\ln{x}$ on $[1, b]$ is approximately $\ln(b) - 1$ when $b$ is large.

\numa Sketch the graph of $f(x) = x e^{-x}$ on the interval $[0, 4]$.

\sub What is the area between the graph of $y = f(x)$ and the $x$-axis
for $0 \leq x \leq 4$?

\sub What is the area between the graph of $y = f(x)$ and the $x$-axis
for $0 \leq x \leq b$?  Express your answer in terms of $b$, and
denote it $A(b)$.  What is $A(100)$?

\num  Find three solutions $y = f(t)$ to the differential equation
	$$
	\dfrac{dy}{dt} = 5 - 2 \ln{t}.
	$$

\numa  Find the solution $y = \varphi(t)$ to the initial value problem
	$$
	\dfrac{dy}{dt} = t e^{-t^2/2},
		\qquad y(0) = 2.
	$$

\sub The function $\varphi(t)$ increases as $t$ increases.  Show this
first by sketching the graph of $y = \varphi(t)$.  Show it also by
referring to the differential equation that $\varphi(t)$ satisfies.
(What is true about the derivative of an increasing function?)

\sub Does the value of $\varphi(t)$ increase without bound as $t \to
\infty$?  If not, what value does $\varphi(t)$ approach?

\numa {\bf The differential form of Integration by Parts}\quad If $u$
and $v$ are expressions in $x$, then the product rule can be written
as
\[
\dfrac{d}{dx}(u\cdot v) = \dfrac{du}{dx}\cdot v + u \cdot \dfrac{dv}{dx}.
\]
Explain carefully how this leads to the following statement of integration by parts, and why it is equivalent to the form in the text:
$$\int u\,dv = uv-\int v\,du.$$

\sub Solve a couple of the preceding problems using this notation.


\subsubsection{Sine and cosine integrals}\index{integral!sine and cosine}

The purpose of the remaining exercises is to establish integral
formulas that we will use to analyze Fourier polynomials and the power
spectrum in chapter~12.  In the first three, $\alpha$ is a constant:
\begin{align*}
\int \sin^2{\alpha x}\,dx &= \dfrac{x}{2} 
	- \dfrac{1}{4\alpha}\sin{2\alpha x} + C,  \\
\int \cos^2{\alpha x}\,dx &= \dfrac{x}{2} 
	+ \dfrac{1}{4\alpha}\sin{2\alpha x} + C,  \\
\int \sin{\alpha x}\cos{\alpha x}\,dx &=  
	- \dfrac{1}{4\alpha}\cos{2\alpha x} + C.
\end{align*}
In the remaining four, $\alpha$ and $\beta$ are {\em different\/} constants:
\begin{align*}
\int \sin{\alpha x} \sin{\beta x}\,dx &= \dfrac{1}{\beta^2-\alpha^2}
	\left(\alpha\cos{\alpha x}\sin{\beta x} - \beta\sin{\alpha x}
	\cos{\beta x}\right) + C, \\
\int \cos{\alpha x} \cos{\beta x}\,dx &= \dfrac{1}{\beta^2-\alpha^2}
	\left(\beta\cos{\alpha x}\sin{\beta x} - \alpha\sin{\alpha x}
	\cos{\beta x}\right) + C, \\
\int \sin{\alpha x} \cos{\beta x}\,dx &= \dfrac{1}{\beta^2-\alpha^2}
	\left(\beta\sin{\alpha x}\sin{\beta x} + \alpha\cos{\alpha x}
	\cos{\beta x}\right) + C, \\
\int \cos{\alpha x} \sin{\beta x}\,dx &=  \dfrac{1}{\beta^2-\alpha^2}
	\left(-\alpha\sin{\alpha x}\sin{\beta x} - \beta\cos{\alpha x}
	\cos{\beta x}\right) + C.
\end{align*}

\numa In the later exercises we shall make frequent use of the
following \label{`trig'} ``trigonometric
identities''\index{trigonometric identities}:
\begin{align*}
    2 \sin{\alpha x} \cos{\alpha x} &= \sin{2\alpha x}, \\
\cos^2{\alpha x} - \sin^2{\alpha x} &=  \cos{2\alpha x}, \\
\sin^2{\alpha x} + \cos^2{\alpha x} &= 1.
\end{align*}
Using a graphing package on a computer, graph together the functions
	$$
    2 \sin{\alpha x} \cos{\alpha x} \qquad \mbox{and} 
		\qquad \sin{2\alpha x} 
	$$
to show that they seem to be identical.  (That is, show that they
``share phosphor.'')  Then do the same for the pairs of functions in
the other two identities.

\sub We can give a different argument for the identities above using
the ideas we have developed in studying initial value problems.  To
prove the first identity, for instance, let $f(x) = 2 \sin{\alpha x}
\cos{\alpha x}$, and let $g(x) = \sin{2\alpha x}$.  Show that both
functions satisfy
$$
y'' = -4\alpha^2 y, \quad\mbox{with}\quad 
    y(0)=0\quad\mbox{and}\quad y'(0)=2\alpha.
$$
Hence conclude the two functions must be the same.

\sub Find an initial value problem that is satisfied by both
$f(x)=\cos^2{\alpha x} - \sin^2{\alpha x}$ and by $g(x) = \cos{2\alpha
  x}$.

\num {\bf Evaluating } $\displaystyle \int \sin^2{x} \, dx$  

\sub  Using integration by parts, show that 
	$$
	\int \sin^2{x} \, dx = 
		- \sin{x} \cos{x} + \int \cos^2{x} \, dx.
	$$

\sub Using the identity $\sin^2{\alpha x} + \cos^2{\alpha x} = 1$,
show that the new integral can be written as
	$$
	x - \int \sin^2{x} \, dx.
	$$

\sub  Combining (a) and (b) algebraically, show that
	$$
	2 \int \sin^2{x} \, dx = - \sin{x} \cos{x} + x + C.
	$$

\sub  Using algebra and a trigonometric identity, conclude that
	$$
	\int \sin^2{x} \, dx = 
		\dfrac{x}{2} - \dfrac{1}{4} \sin{2x} + C.
	$$

\num  Modify the argument of the preceding exercise to show		%qqq
	$$
	\int \sin^2{\alpha x}\,dx  =  
	  \dfrac{x}{2} - \dfrac{1}{4\alpha}\sin{2\alpha x} + C.
	$$
and
	$$
	\int \cos^2{\alpha x}\,dx  =  
	  \dfrac{x}{2} + \dfrac{1}{4\alpha}\sin{2\alpha x} + C.
	$$


\num {\bf Evaluating }$\displaystyle \int \sin^2{\alpha x} \, dx$  
\smallskip

\noindent
Determine this integral anew, without using integration by parts, by
carrying out the following steps.

\sub From the trigonometric identities\index{trigonometric identities}
on page~\pageref{`trig'}, deduce that
	$$
	2 \sin^2 {\alpha x} = 1 - \cos{2 \alpha x}.
	$$

\sub  Using the formula in (a), conclude that
	$$
	\int \sin^2{\alpha x} \, dx = 
	\dfrac{x}{2} - \dfrac{1}{4\alpha}\sin{2\alpha x} + C.
	$$

\num {\bf Evaluating }$\displaystyle \int \cos^2{\alpha x} \, dx$  
\smallskip

\noindent
Using only algebra and the identity $\sin^2{\alpha x} + \cos^2{\alpha
  x} = 1$, show that the previous exercise implies
	$$
	\int \cos^2{\alpha x} \, dx = 
	\dfrac{x}{2} + \dfrac{1}{4\alpha}\sin{2\alpha x} + C.
	$$

\num {\bf Evaluating }$\displaystyle 
		\int \sin{\alpha x} \cos{\alpha x} \, dx$  

\sub Using the identity $\sin{2\alpha x} = 2 \sin{\alpha x}
\cos{\alpha x}$, deduce the following formula
	$$
	\int \sin{\alpha x} \cos{\alpha x} \, dx = 
		-\dfrac{1}{4\alpha} \cos{ 2 \alpha x} + C.
	$$

\sub  Using integration by substitution, obtain the alternative formula
	$$
	\int \sin{\alpha x} \cos{\alpha x} \, dx = 
		\dfrac{1}{2\alpha} \sin^2{\alpha x} + C.
	$$

\sub Show that your results in (a) and (b) are compatible.  (For
example, use exercise 15~(a).)

\num 
{\bf Evaluating }${\displaystyle \int \sin{\alpha x} \sin{\beta x}\, dx}$


\sub  Use integration by parts to show that 
	$$
	\int \sin{\alpha x} \sin{\beta x}\, dx =
		-\dfrac{1}{\alpha} \cos{\alpha x} \sin{\beta x}
		+ \dfrac{\beta}{\alpha} 
		\int \cos{\alpha x} \cos{\beta x}\, dx.
	$$

\sub Using integration by parts again, show that the new integral in
part (a) can be written as
	$$
	\int \cos{\alpha x} \cos{\beta x}\, dx =
		\dfrac{1}{\alpha} \sin{\alpha x} \cos{\beta x}
		+ \dfrac{\beta}{\alpha} 
		\int \sin{\alpha x} \sin{\beta x}\, dx.
	$$

\sub Let $J = \dint \sin{\alpha x} \sin{\beta x}\, dx$; show that
combining (a) and (b) gives
	$$
	J = -\dfrac{1}{\alpha} \cos{\alpha x} \sin{\beta x}
	+ \dfrac{\beta}{\alpha^2} \sin{\alpha x} \cos{\beta x}
	+ \dfrac{\beta^2}{\alpha^2} J.
	$$

\sub  Solve (c) for $J$ to find
	$$
	\int \sin{\alpha x} \sin{\beta x}\, dx = 
	 \dfrac{1}{\beta^2 - \alpha^2}
		(\alpha \cos{\alpha x} \sin{\beta x} 
		- \beta \sin{\alpha x} \cos{\beta x}) + C.
	$$

\num 	Imitate the methods of the preceding exercise to deduce		%qqq
	$$
	\int \cos{\alpha x} \cos{\beta x}\, dx =
	\dfrac{1}{\beta^2 - \alpha^2} 
		(\beta \cos{\alpha x} \sin{\beta x} 
		- \alpha \sin{\alpha x} \cos{\beta x}) + C
	$$
and
	$$
	\int \sin{\alpha x} \cos{\beta x}\, dx =
	\dfrac{1}{\beta^2 - \alpha^2}
		(\beta \sin{\alpha x} \sin{\beta x} 
		+ \alpha \cos{\alpha x} \cos{\beta x}) + C.
	$$

\num 
{\bf Evaluating }${\displaystyle \int \cos{\alpha x} \sin{\beta x}\, dx}$

\smallskip

This integral is the same as one in the preceding exercise, if you exchange the factors $\alpha$ and $\beta$.  Do that, and obtain the formula  %qqq
	$$
	\int \cos{\alpha x} \sin{\beta x}\, dx =
	{1}{\alpha^2 - \beta^2} 
		(\alpha \sin{\alpha x} \sin{\beta x} 
		+ \beta \cos{\alpha x} \cos{\beta x}) +C.
	$$

\vspace{-10pt}

\num  Determine the following.

\noindent
\begin{minipage}{195pt}
\sub ${\displaystyle \int_0^{2\pi} \sin^2{x}\, dx}$.

\sub ${\displaystyle \int_0^{n\pi} \cos^2{x}\, dx}$, $n$ a positive integer.
\end{minipage}
\hfill
\begin{minipage}{155pt}
\sub ${\displaystyle \int_0^{2\pi} \sin{x} 
\sin{\beta x}\, dx}$, $\beta \neq 1$.

\sub ${\displaystyle \int_0^{\pi} \sin{x} 
\sin{\beta x}\, dx}$, $\beta \neq 1$.
\end{minipage}
\hfill
\mbox{}

%%%%%%%%%%%%%%%%%%%%%%%%%%%%%%%%%%%%%%%%%%%%%%%%%%%%%%%%%%%%%%%%%%%%%%%%
%                                                                      %
%                               Section 4                              %
%                                                                      %
%%%%%%%%%%%%%%%%%%%%%%%%%%%%%%%%%%%%%%%%%%%%%%%%%%%%%%%%%%%%%%%%%%%%%%%%

\section{Separation of Variables and\\ Partial Fractions}

One of the principal uses of integration techniques is to find closed
form solutions to differential equations.  If you look back at the
methods we have developed so far in this chapter, they are all
applicable to differential equations of the form $y'=f(t)$ for some
function $f$---that is, the rate at which $y$ changes is a function of
the independent variable only.  In such cases we only need to find an
antiderivative $F$ for $f$, choose the constant $C$ to satisfy the
initial value, and we have our solution.  As we saw in the early
chapters, though, the behavior of $y'$ often depends on the values of
$y$ rather than on $t$---think of the $S$-$I$-$R$ model or the various
predator-prey problems.  In this section we will see how our earlier
techniques can be adapted to apply to problems of this sort as well.

\subsection{The Differential Equation $y' = y$}

As you know, 
%
\mar{A new method for solving $y' = y$}	
%
the exponential functions $y = C e^t$ are the solutions to the
differential equation $dy/dt = y$.  Let's put aside this knowledge for
a moment and rediscover these solutions using a new method.  The
method involves the connection between inverse functions and their
derivatives.  With it we will be able to explore a variety of problems
that had been beyond our reach.

The idea behind the new method is quite simple:  
%
\mar{Find the inverse function instead}
%
Instead of thinking of $y$ as a function of $t$, convert to thinking
of $t$ as a function of $y$, thereby looking for the {\bf inverse
  function}.  We know that the derivative of the inverse function is
the reciprocal of the derivative of the original function, so we can
rewrite the given differential equation by using its reciprocal:
\mar{\vspace{12pt} The new differential equation\ldots}
	$$
	\dfrac{dy}{dt} = y \qquad \mbox{becomes} \qquad 
		\dfrac{dt}{dy} = \dfrac{1}{y}.
	$$
Then solve the new differential equation $dt/dy = 1/y$.  While this
may not look very different, it has the property that the rate of
change of the {\em dependent variable}---now $t$---is expressed as a
function of the {\em independent} variable---now $y$.  But this is
just the form we have been considering in the earlier sections of this
chapter.  A solution to the new equation is a function $t = g(y)$
whose derivative is $1/y$.  This is one of the basic antiderivatives
listed 
%
\mar{\vspace{12pt} \ldots and its solution}
%
in the table on page~\pageref{table}:
	$$
	t = g(y) = \ln{y} + k,
	$$
where $k$ is an arbitrary constant.  

The solution to the original differential equation 
%
\mar{The inverse \ldots}	
%
$dy/dt = y$ is the inverse of $t = \ln{y} + k$.  We find it by solving
this equation for $y$:
\begin{align*}
t - k &= \ln{y} \\
e^{t - k} &= y \\
e^t \cdot e^{-k} &= y
\end{align*}
Thus $y(t) = Ce^t$, 
%
\mar{\vspace{-18pt} \ldots solves the\\ original problem}
%
where we have replaced the constant $e^{-k}$ by $C$.
\bigskip

\subsubsection{Indefinite integrals}
\index{indefinite integral}

The language of indefinite integrals 
%
\mar{The differential equation expressed using differentials}
%
and differentials again provides a convenient mnemonic for this new
method.  First, we use the original differential equation to relate
the differentials $dy$ and $dt$:
	$$
	dy = \dfrac{dy}{dt} \, dt = y \, dt.
	$$
This use of differentials is introduced on
page~\pageref{differentials}.  At this point the equation makes sense
for either $t$ or $y$ being the independent variable.  Now if we try
to integrate this equation with respect to $t$, we get
	$$
	y = \int dy = \int y(t) \, dt.
	$$
We can't find the last integral, because we don't know what $y$ is as
a function of $t$.  Remember, an indefinite integral is an
antiderivative, so an expression of the form
	$$
	\int f \, dt
	$$
represents a function $F(t)$ whose derivative is $f(t)$.  If $f$ is
{\em not\/} given by a formula in $t$, there is no way to get a
formula for $F(t)$.

Suppose, though, that we divide both sides of the differential
equation $dy = y \, dt$ by $y$, and integrate with respect to $y$.
The equation takes the form
	$$
	\dfrac{dy}{y} = dt.	
	$$
Now, if we introduce indefinite integrals, we have
	$$
	\int \dfrac{dy}{y} = \int dt.
	$$
This time the variables $y$ and $t$ 
%
\mar{The variables\\ are now \emph{separated}}
%
have been separated from each other, and we {\em can\/} find the
integrals.  In fact,
	$$
	\ln{y} = \int \dfrac{dy}{y} = \int dt = t + b,
	$$
where $b$ is an arbitrary constant. (We could just as easily have
added the constant to the left side instead---do you see why we don't
have to add a constant to both sides?)  To complete the work 
%
\mar{\vspace{10pt}The solution\\ once again}
%
we solve for~$y$:
	$$
	y = e^{t + b} = e^t \cdot e^b = Ce^t,
	$$
where $C = e^b$.


\subsubsection{Summary}

The first time we went through the method, we replaced 
	$$
\dfrac{dy}{dt} = y \qquad \mbox{by} \qquad \dfrac{dt}{dy} = \dfrac{1}{y}.
	$$
These differential equations express the same relation between $y$ and
$t$.  Each is just the reciprocal of the other.  In the first, $y$
depends on $t$; in the second, though, $t$ depends on $y$.  The second
time we went through the method, using indefinite integrals, we
replaced
	$$
	dy = y \, dt \qquad \mbox{by} \qquad dt = \dfrac{dy}{y}.
	$$
This change was also algebraic, and it had the same effect: the
dependent variable changed from $y$ to $t$.  More important,
%
\mar{To integrate,\\ separate the variables}
%
in the new differential equation (using the differentials themselves!)
the variables are separated.  That allows us to do the integration.
We get a solution in the form $t = g(y)$.  The solution to the
original problem is the {\em inverse\/} $y = f(t)$ of $t = g(y)$.

\subsection{Separation of Variables}
\index{separation of variables}

With the method of {\bf separation of variables}, introduced in the
previous pages, we can obtain formulas for solutions to a number of
differential equations that were previously accessible only by Euler's
method.  Recall that one of the clear advantages of a {\em formula\/}
is that it allows us to see how the parameters in the problem affect
the solution.  We'll look at two problems.  First we'll show how the
method can explain the rather baffling formula for supergrowth that we
gave in chapter~4.  Then, using the method of {\bf partial fractions},
to be discussed next, we'll give a formula for logistic growth.
	
\subsubsection{Supergrowth}

In chapter 4.2 we modelled the growth of a population $Q$ by the
initial value problem\index{supergrowth}
	$$
	\dfrac{dQ}{dt} = k Q^{1.2}, \qquad Q(0) = A.
	$$
To get a formula for the solution, 
%
\mar{Separate the variables}
%
transform the differential equation in the following way:
	$$
	\dfrac{dQ}{dt} = k \, Q^{1.2} \quad \leadsto \quad 
		dQ = k \, Q^{1.2} \, dt \quad \leadsto \quad
		\dfrac{dQ}{Q^{1.2}} = k \, dt.
	$$
Now integrate:
	$$
	\int \dfrac{dQ}{Q^{1.2}} = \int k \, dt.
	$$
Because the variables have been separated, the integrals can be found:
\begin{align*}
\dint \dfrac{dQ}{Q^{1.2}} 
   &= \int Q^{-1.2} \, dQ = \dfrac{1}{-.2} Q^{-.2} \\[3pt]
\dint k \, dt &= kt + C
\end{align*}
Therefore, $\tfrac{1}{-.2} Q^{-.2} = kt + C$.  
%
\mar{Solve for $Q$}
%
Now we must solve this equation for $Q$.  Here is one possible
approach.  First, we can write
	$$
	Q^{-.2} = -.2(kt + C) \ = \ C_1 -.2kt.
	$$
To simplify the expression, we have replaced $-.2C$ by a {\em new\/}
constant $C_1$.  Since
	$$
	(Q^{-.2})^{-5} = Q^{-.2 \times -5} = Q^1 = Q,
	$$ 
we'll raise both sides of the previous equation to the power $-5$:
	$$
	Q(t) = (Q^{-.2})^{-5} = (C_1 -.2kt)^{-5}.
	$$
	
The last step is to incorporate 
%
\mar{Bring in the\\ initial condition}
%
the initial condition $Q(0) = A$.  According to the new formula for
$Q(t)$,
	$$
	Q(0) = (C_1 -.2k\cdot 0)^{-5} = C_1^{-5} = A.
	$$
Solving $A = C_1^{-5}$ for $C_1$, we get 
	$$
	C_1 = A^{-1/5} = \dfrac{1}{\sqrt[5]{A}}.
	$$
We are now done:
	$$
	Q(t) = \left(\dfrac{1}{\sqrt[5]{A}} -.2kt \right)^{-5}.
	$$
This is the formula that appears on page~\pageref{supergrowth soln}.  
%
\mar{Interpreting\\ the formula}
%
It shows how the parameters $k$ and $A$ affect the solution.  
In particular, we called this {\em supergrowth\/} because the model predicts that the population $Q$ becomes infinite when 
	$$
	\dfrac{1}{\sqrt[5]{A}} -.2kt = 0; \quad \mbox{that is, when}
		\quad t = \dfrac{1}{.2k \sqrt[5]{A}}.
	$$ 	

\subsection{Partial Fractions}
\index{partial fractions}

Using separation of variables with a partial fractions decomposition
(to be described below), we will obtain a formula for the solution to
the logistic equation (chapter~4.1).  The method of {\bf partial
  fractions} is a useful tool for solving many integration problems.

\subsubsection{Logistic growth}
\index{logistic growth}	

Consider this initial value problem associated with the logistic differential equation:
$$
\dfrac{dP}{dt} = kP \left( 1 - \dfrac{P}{C} \right), \qquad P(0) = A.
$$
We will find a formula for the solution that incorporates the growth
parameter $k$ and the carrying capacity $C$.

The first step is to transform the equation 
%
\mar{Step 1: separate\\ the variables}	
%
into one where the variables are separated:
	$$
	\dfrac{dP}{dt} = kP \left( 1 - \dfrac{P}{C} \right) \quad \leadsto
		\quad \dfrac{dP}{P (1 - P/C)} = k \, dt.
	$$
Integrating the new equation we get
	$$
	\int \dfrac{dP}{P (1 - P/C)} = \int k \, dt.
	$$
We are stuck now, 
%
\mar{The denominator has an unfamiliar form}
%
because the integral on the left doesn't appear in our table of
integrals (page~\pageref{table}).  If the denominator had {\em only\/}
$P$ or {\em only\/} $1 - P/C$, we could use the natural logarithm.
The difficulty is that the denominator is the product of both terms.

There is a way out of the difficulty.  We will use algebra to
transform the integrand into a form we can work with.  The first step
is to simplify the denominator a bit:
\[
\dfrac{1}{P(1 - P/C)} = \dfrac{1}{P(C/C - P/C)} = \dfrac{1}{P(C - P)/C}
		= \dfrac{C}{P(C - P)}.
\]
(This wasn't essential; it just makes later steps easier to write.)
The next step will be the crucial one.  To understand why we take it,
consider the rule for adding two fractions:
	$$
	\dfrac{\alpha}{(x + a)} + \dfrac{\beta}{(x + b)} = 
		\dfrac{\alpha (x + b) + \beta (x + a)}{(x + a) (x + b)}
	$$
The denominator is a product---very much like the product $P(C - P)$ 
%
\mar{Step 2: write the integrand as a sum\\ of simple fractions}
%
in our integrand!  Perhaps we can write {\em that\/} as a sum of two
simpler fractions:
	$$
\dfrac{C}{P(C - P)} = \dfrac{\alpha}{P} + \dfrac{\beta}{C - P}.
	$$
What values should $\alpha$ and $\beta$ have?  According to the rule
for adding fractions,
	$$
	\dfrac{\alpha}{P} + \dfrac{\beta}{C - P} = 
		\dfrac{\alpha (C - P) + \beta P}{P(C - P)},
	$$
and this should equal the original integrand:
	$$
	\dfrac{\alpha (C - P) + \beta P}{P(C - P)} = \dfrac{C}{P(C - P)}.
	$$

\noindent
Since the denominators are equal, 
%
\mar{Determining \\ $\alpha$ and $\beta$}
%
the numerators must also be equal:
	$$
	\alpha (C - P) + \beta P = C.
	$$
In fact, they must be equal {\em as polynomials in the variable\/} $P$.  If we rewrite the last equation, collecting terms that involve the same power of $P$, we get 	$$
	(\beta - \alpha)P + \alpha C = 0 \cdot P + 1 \cdot C.
	$$
Since two polynomials are equal precisely when their coefficients are
equal, it follows that
	$$
	\beta - \alpha = 0 \qquad \alpha = 1.
	$$
Thus $\alpha = \beta = 1$, 
%
\mar{\vspace{16pt} The partial fractions decomposition}
%
and we have 
	$$
	\dfrac{C}{P(C - P)} = \dfrac{1}{P} + \dfrac{1}{C - P}.
	$$
The simpler expressions on the right are called {\bf partial
  fractions}.\index{partial fractions} Their denominators are the
different {\em parts\/} of the denominator of the integrand.  The
equation that expresses the integrand as a sum of partial fractions is
called a {\bf partial fractions decomposition}.\index{partial
  fractions decomposition}
\medskip

We can now return to the integral equation we are trying to solve:
	$$
	\int \dfrac{C}{P(C - P)} \, dP  = \int k \, dt.
	$$

\newpage

\noindent
The right hand side equals $kt	+ b$, 
%
\mar{Step 3: evaluate\\ the integrals}
%
where $b$ is the usual constant of integration.  Thanks to the partial
fractions decomposition, the left hand side can be written
	$$	
	\int \dfrac{C}{P(C - P)} \, dP  = \int \dfrac{1}{P} \, dP + 
		\int \dfrac{1}{C - P} \, dP.
	$$
The first integral on the right is straightforward:
	$$
	\int \dfrac{1}{P} \, dP = \ln{P}.
	$$
The second can be solved by using the substitution $C - P = u$, with
$dP = -du$:
	$$
	\int \dfrac{1}{C - P} \, dP = \int \dfrac{-du}{u} = 
		- \ln{u} = -\ln(C - P).
	$$
Putting everything together we find
	$$
	\ln{P} - \ln(C - P) = \ln \left( \dfrac{P}{C - P} \right) = kt + b.
	$$

As we have seen, 
%
\mar{Step 4:\\ solve for $P$}	
%
separation of variables usually leaves us with an inverse function to
find.  This problem is no different.  We must solve the last equation
for $P$.  The first step is to exponentiate both sides:
	$$
	\dfrac{P}{C - P} = e^{kt + b} = e^b \cdot e^{kt} = Be^{kt}.
	$$
To simplify the expression a bit, we have replaced $e^b$ by $B = e^b$.
Multiplying both sides by $C - P$ gives
	$$
	P = Be^{kt}(C - P) = CBe^{kt} - PBe^{kt}.
	$$
Now bring the last term over to the left, and then factor out $P$:
	$$
	P + PBe^{kt} = \left( 1 + Be^{kt} \right) P = CBe^{kt}.
	$$
The final step is to divide 
%
\mar{The formula for $P(t)$}
%
by the coefficient $1 + Be^{kt}$:
	$$
	P(t) = \dfrac{CBe^{kt}}{1 + Be^{kt}}.
	$$


Lastly, 
%
\mar{Step 5: incorporate\\ the initial condition}
%
we must see how the initial condition $P(0) = A$ affects the solution.
We could substitute $t = 0$, $P = A$ into the last formula, but that
produces an algebraic mess.  
%
\pagebreak
%
We want to know how $A$ affects the constant $B$, and we can see that
directly by making our substitutions into the equation above:
	$$
	\dfrac{P}{C - P} = Be^{kt} \quad \leadsto \quad
		\dfrac{A}{C - A} = Be^0 = B.
	$$
Now replace $B$ by $A/(C - A)$ in our formula for $P(t)$.  This yields
	$$
	P(t) = \dfrac{CAe^{kt}/(C - A)}{1 + Ae^{kt}/(C - A)} =
		\dfrac{CAe^{kt}}{C - A + Ae^{kt}}
	$$
If we write $-A + Ae^{kt} = A(e^{kt} - 1)$, we get one of the standard
forms of the solution to 
%
\mar{\vspace{28pt} The complete solution}
%
the logistic equation: 
	$$
	P(t) = \dfrac{CAe^{kt}}{C + A(e^{kt} - 1)}.
	$$

\noindent
{\bf Remark}.  The method of partial fractions\index{partial
  fractions} can be used to evaluate integrals of the form
	$$
	\int \dfrac{dx}{(x + a_1) (x + a_2) \cdots (x + a_n)} \qquad 
	\mbox{or} \qquad \int \dfrac{P(x)}{Q(x)} \, dx,
	$$
where $P(x)$ and $Q(x)$ are arbitrary polynomials.  Tables of
integrals and calculus references describe how the method works in
these cases.  As one example, though, let's compute the antiderivative
of the cosecant function, since we will need it in the next section.

\medskip
\noindent
{\bf Example---The antiderivative of the
  cosecant}\quad\label{trick}\index{cosecant!antiderivative of} We
first use a trigonometric identity to transform the integral slightly:
\begin{align*}
\dint \csc{x}\,dx &= \dint \dfrac{1}{\sin{x}}\,dx \\[6pt]
  &= \dint \dfrac{\sin{x}}{\sin^2{x}}\,dx \\[6pt]
  &= \dint \dfrac{\sin{x}\,dx}{1-\cos^2{x}}.
\end{align*}
If we now make the substitution $u=\cos{x}$, with $du=-\sin{x}\,dx$,
this becomes
\begin{align*}
\dint \csc{x}\,dx &= \dint \dfrac{-du}{1-u^2} \\[3pt]
 &= \dint \dfrac{-1}{2}\left(\dfrac{1}{1+u} 
            + \dfrac{1}{1-u}\right)\,du \\[3pt]
 &= -\dfrac{1}{2}\dint  \dfrac{du}{1+u} 
            - \dfrac{1}{2}\dint \dfrac{du}{1-u} \\[3pt]
 &= -\dfrac{1}{2}\ln(1+u) + \dfrac{1}{2} \ln(1-u) +C\\[6pt]
 &= \dfrac{1}{2}\ln{\dfrac{1-u}{1+u}} +C\\[6pt]
 &= \dfrac{1}{2}\ln{\dfrac{1-\cos{x}}{1+\cos{x}}}+C.
\end{align*}
We can simplify this slightly 
%
\mar{The final form\\ of the antiderivative\\ of the cosecant}
%
by multiplying both numerator and denominator by $1+\cos{x}$ to get
\begin{align*}
\dint \csc{x}\,dx
  &= \dfrac{1}{2}\ln\dfrac{1-\cos^2{x}}{(1+\cos{x})^2}+C
   = \ln{\left|\dfrac{\sin{x}}{1+\cos{x}}\right|}+C \\
  &=-\ln{|\csc{x}+\cot{x}|}+C.\label{cosecant}
\end{align*}
Note that in the next to last line we used the general fact about
logarithms that $n\,\ln{A}=\ln(A^n)$ for any value of $n$ and any
$A>0$.  Note also that since the domain of the secant function
consists of infinitely many separate intervals. the ``$+\ C$'' at
the end of the antiderivative needs to be interpreted as potentially a
different value of $C$ over each interval.

In the same fashion we can 
%
\mar{The antiderivative\\ of the secant}
%
obtain the antiderivative for the secant function\index{secant!antiderivative of}:\label{secant}
$$
\int \sec{x}\,dx=\ln{|\sec{x} +\tan{x}|}+C.
$$


\subsection{Exercises}

\subsubsection{Separation of variables}
\index{separation of variables}

\vspace{-10pt}

\num Use the method of separation of variables to find a formula for
the solution of the differential equation $dy/dt = y + 5$.  Your
formula should contain an arbitrary constant to reflect the fact that
many functions solve the differential equation.

\newpage

\num Use the method of separation of variables to find formulas for
the solutions to the following differential equations.  In each case
your formula should be expressed in terms of the input variable that
is indicated (e.g., in part (a) it is $t$).

\noindent
\begin{minipage}[t]{150pt}
\sub  $dy/dt = 1/y$.

\sub  $dz/dx = 3/(z - 2)$.

\sub  $dy/dx = x/y$
\end{minipage}
\hfill
\begin{minipage}[t]{150pt}
\sub  $dy/dx = y/x$

\sub  $du/dv = u/(u - 1)$

\sub  $dv/dt = -\sqrt{v}$
\end{minipage}
\hfill
\mbox{}

\num {\bf A cooling liquid}.  According to Newton's law of
cooling\index{Newton's law of cooling} (see chapter~4.1), in a room
where the ambient temperature is $A$, the temperature $Q$ of a hot
object will change according to the differential equation
	$$
	\dfrac{dQ}{dt} = -k(Q - A).
	$$
The constant $k$ gives the rate at which the object cools.

\sub Find a formula for the solution to this equation using the method
of separation of variables.  Your formula should contain an arbitrary
constant.

\sub Suppose $A$ is 20$^\circ$C and $k$ is .1$^\circ$ per minute per
$^\circ$C.  If time $t$ is measured in minutes, and $Q(0) =
90^\circ$C, what will $Q$ be after 20 minutes?

\sub  How long does it take for the temperature to drop to 30$^\circ$C?

\numa Suppose a cold drink at 36$^\circ$F is sitting in the open air
on a summer day when the temperature is 90$^\circ$F.  If the drink
warms up at a rate of .2$^\circ$F per minute per $^\circ$F of
temperature difference, write a differential equation to model what
will happen to the temperature of the drink over time.

\sub Obtain a formula for the temperature of the drink as a function
of the number of minutes $t$ that have passed since its temperature
was 36$^\circ$F.

\sub What will the temperature of the drink be after 5 minutes; after
10 minutes?

\sub  How long will it take for the drink to reach 55$^\circ$F?

\num {\bf A leaking tank}.\index{leaking tank} In chapter~4.2 we used
the differential equation
	$$
	\dfrac{dV}{dt} = -k \sqrt{V}
	$$
to model the volume $V(t)$ of water in a leaking tank after $t$ hours
(see page~\pageref{leaking tank-begin}).

\sub  Use the method of separation of variables to show that 
	$$
	V(t) = \dfrac{k^2}{4} \, (C - t)^2
	$$
is a solution to the differential equation, for any value of the
constant~$C$.

\sub  Explain why the function
\[
V(t) = \begin{cases}
\dfrac{{k^2}}{4} (C - t)^2 &\text{if $0 \leq t \leq C$}, \\[8pt]
0 &  \text{if $C < t$}.
\end{cases}
\]
is {\em also\/} a solution to the differential equation.  Why is {\em
  this\/} solution more relevant to the leaking tank problem than the
solution in part~(a)?

\num {\bf A falling body with air resistance}.\index{falling body} We
have used the differential equation
	$$
	\dfrac{dv}{dt} = -g -bv
	$$ to model the motion of a body falling under the influence
        of gravity ($g$) and air resistance ($bv$).  Here $v$ is the
        velocity of the body at time $t$.  (See pages \pageref{fall
          w/air-begin}--\pageref{fall w/air-end}.)

\sub  Solve the differential equation by separating variables, and obtain
	$$
	v(t) = \dfrac{1}{b} \, \left( Ce^{-bt} - g \right),
	$$
where $C$ is an arbitrary constant.

\sub Now impose the initial condition $v(0) = 0$ (so the body starts
it fall from rest) to determine the value of $C$.  What is the formula
for $v(t)$ now?

\sub Exercise 21 on page~\pageref{ex:falling body1} gives the solution
to the initial value problem as
	$$							
	v(t) = \dfrac{g}{b} \, \left( 2^{-bt/.69} - 1 \right).	
	$$
Reconcile this expression with the one you obtained in part (b) of
this exercise.

\sub The distance $x(t)$ that the body has fallen by time $t$ is given
by the integral
	$$
	x(t) = \int_0^t v(t) \, dt, \quad \mbox{because} \quad
		\dfrac{dx}{dt} = v \quad \mbox{and} \quad x(0) = 0.
	$$
Use your formula for $v(t)$ from part (b) to find $x(t)$.

\newpage

\numa {\bf Supergrowth}.\index{supergrowth} We have analyzed the
differential equation
	$$
	\dfrac{dQ}{dt} = k Q^p
	$$
when $p = 1.2$ (and, of course, when $p = 1$).  Find a formula for the
solution $Q(t)$ when $p = 2$.  Your formula should contain an
arbitrary constant $C$.
  
\sub Add the initial condition $Q(0) = A$.  This will fix the value of
the constant~$C$.  What is the formula for $Q(t)$ when the initial
condition is incorporated?

\sub Your formula in part (b) should demonstrate that $Q$ becomes
infinite at some finite time $t = \tau$.  When is $\tau$?  Your answer
should be expressed in terms of the growth constant $k$ and the
initial population size~$A$.

\sub Suppose the values of $k$ and $A$ are known only imprecisely, and
they could be in error by as much as 5\%.  That makes the value of
$\tau$ uncertain.  Which error causes the greater uncertainty: the
error in $k$ or the error in $A$?  (See the discussion of error
analysis for the supergrowth model on pages \pageref{supergrowth
  errors-begin}--\pageref{supergrowth errors-end}.)

\num {\bf General supergrowth}.  Find the solution to the initial
value problem
	$$
	\dfrac{dQ}{dt} = k Q^p, \qquad Q(0) = A
	$$
for {\em any\/} value of the power $p$.  For which values of $p$ does
$Q$ blow up to $\infty$ at a finite time $t = \tau$?  What is $\tau$?



\subsubsection{Partial fractions}
\index{partial fractions}

\num Use the method of partial fractions to determine the values of
$\alpha$, $\beta$, and $\gamma$ in the following equations.

\noindent
\begin{minipage}{180pt}
\sub $\dfrac{1}{(x - 1)(x + 2)} = \dfrac{\alpha}{x - 1} +
	\dfrac{\beta}{x + 2}$
	
\sub $\dfrac{x}{(x - 1)(x + 2)} = \dfrac{\alpha}{x - 1} +
	\dfrac{\beta}{x + 2}$

\sub $\dfrac{1}{x(x^2-1)} = \dfrac{\alpha}{x} + 
	\dfrac{\beta}{x-1} + \dfrac{\gamma}{x+1}$
\end{minipage}
\hfill
\begin{minipage}{180pt}
\sub $\dfrac{x}{2x^2 +3x + 1} = \dfrac{\alpha}{2x + 1} + 
	\dfrac{\beta}{x + 1}$

\sub  $\dfrac{1}{x(x^2 + 1)} = \dfrac{\alpha}{x} +
	\dfrac{\beta x + \gamma}{x^2 + 1}$

\bigskip

\noindent
[Note that $x^2 + 1$ can't be factored.]
\vfill
\mbox{}

\end{minipage}
\hfill
\mbox{}

\pagebreak

\num  Find a formula for each of these indefinite integrals.
\vspace{-5pt}

\noindent
\begin{minipage}[t]{150pt}
\sub  $\dint \dfrac{3 \, dx}{(x - 1)(x + 2)}$

\sub  $\dint \dfrac{5x + 3}{(x - 1)(x + 2)} \, dx$

\sub  $\dint \dfrac{dt}{t(t^2-1)}$
\end{minipage}
\hfill
\begin{minipage}[t]{150pt}
\sub  $\dint \dfrac{x \, dx}{1 - x^2}$

\sub  $\dint \dfrac{1 - u}{u^2 - 4} \, du$

\sub  $\dint \dfrac{x^2 + 2x + 1}{x(x^2 + 1)} \, dx$
\end{minipage}
\hfill
\mbox{}

\num  Determine
\vspace{-5pt}

\noindent
\begin{minipage}[t]{150pt}
\sub  $\dint_2^3 \dfrac{3 \, dx}{(x - 1)(x + 2)}$

\sub  $\dint_2^4 \dfrac{dt}{t(t^2-1)}$
\end{minipage}
\hfill
\begin{minipage}[t]{150pt}
\sub  $\dint_0^{\pi/4} \dfrac{x \, dx}{1 - x^2}$

\sub  $\dint_1^{\sqrt{3}} \dfrac{x^2 + 2x + 1}{x(x^2 + 1)} \, dx$
\end{minipage}
\hfill
\mbox{}

\num Mirror the derivation of $\dint \csc{x}\,dx$ to find $\dint
\sec{x}\,dx.$

\num  Consider the particular logistic growth model defined by
	$$
\dfrac{dP}{dt} = .2 P\left( 1 - \dfrac{P}{10} \right) \ \mbox{lbs/hr},
	 \qquad P(0) = .5 \ \mbox{lbs}
	$$ (Compare this with the fermentation problems, pages
         \pageref{ferment-begin}--\pageref{ferment-end}.)

\sub  Obtain the formula for the solution to this initial value problem.

\sub  How large will $P$ be after 3 hours; after 10 hours?

\sub When will $P$ reach one-half the carrying capacity--that is, for
which $t$ is $P = 5$ lbs?

\num Derive the formula for $\dint \sec{x}\,dx$ given on
page~pageref{secant}, using methods similar to those used to find an
antiderivative for the cosecant function.


\newpage
%%%%%%%%%%%%%%%%%%%%%%%%%%%%%%%%%%%%%%%%%%%%%%%%%%%%%%%%%%%%%%%%%%%%%%%%
%                                                                      %
%                               Section 5                              %
%                                                                      %
%%%%%%%%%%%%%%%%%%%%%%%%%%%%%%%%%%%%%%%%%%%%%%%%%%%%%%%%%%%%%%%%%%%%%%%%

\section{Trigonometric Integrals}

The preceding sections have covered the main integration techniques
and concepts likely to be needed by most users of calculus.  These
techniques, together with the numerical methods discussed in
chapter~11.6, should be part of the basic tool kit of every
practitioner of calculus.  For those going on in physics or
mathematics, there are additional methods, largely involving
trigonometric functions in various ways, that are sometimes useful.
The purpose of this section is to develop the most commonly used of
these techniques.

Recall that there are only a few simple antiderivatives we can write
down immediately by inspection.  All non-numerical integration
techniques consist of finding transformations that will reduce some
new class of integration problems to a class we already know how to
solve.  Once we have a new class of solvable problems, then we look
for other classes of problems that can be reduced to this new class,
and so on. The techniques we will be developing in this section
involve ways of making such transformations through the use of basic
trigonometric identities, typically in conjunction with integration by
parts or by substitution.  Before we proceed with the integration
techniques, it will be helpful to list the trigonometric identities
used.

\subsubsection{Review of trigonometric identities}

The most frequently used identity is
$$
\sin^2{x} + \cos^2{x} =1,
$$
and the equivalent form obtained by dividing through by $\cos^2{x}$:
$$
\tan^2{x}+1=\sec^2{x},
$$
and by $\sin^2{x}$:
$$
1+\cot^2{x}=\csc^2{x}.
$$ 
The only other identities you will need have already been encountered:
$$
\sin{2x}=2\sin{x}\cos{x}\quad \mbox{and}\quad 
     \cos{2x}=\cos^2{x}-\sin^2{x},
$$
plus the two other forms of the second of these identities,
$$\cos^2{x}=\dfrac{1}{2}(1+\cos{2x})\quad \mbox{and}\quad
     \sin^2{x} =\dfrac{1}{2}(1-\cos{2x}).
$$

\subsection{Inverse Substitution}\index{inverse substitution!integration by}\index{integration by inverse substitution}

The method of substitution outlined in chapter~11.2 worked by taking a
complicated integrand and breaking it down into simpler components,
reducing the problem 
%
\mar{Success sometimes comes by making things more complicated} 
%
of finding an antiderivative for something in the
form $f(G(x))g(x)$ to the problem of finding an antiderivative for
$f$.  In some cases, though, we go in the opposite direction: we have
an integral $\int f(x)\,dx$ we want to find but can't evaluate
directly.  Instead, we can find a function $G(u)$ with derivative
$g(u)$ such that we can find an antiderivative for $f(G(u))g(u)$.
Since we know this integral is $F(G(u))$, we can now figure out what
the desired function $F$ must be.  As with the earlier substitution
techniques, this {\bf inverse substitution} is conveniently expressed
using differential notation.

\medskip
\noindent
{\bf Example 1}\quad  Suppose we want to evaluate
$$\int \sqrt{4-x^2}\,dx.$$
If we substitute $x=2\sin{u}$, so that $dx=2\cos{u}\,du$, look what happens:
\begin{align*}
\int \sqrt{4-x^2}\,dx
        &= \int \sqrt{4-(2\sin{u})^2}\,2\cos{u}\,du \\
	&= \int \sqrt{4-4\sin^2{u}}\,2\cos{u}\,du \\
	&= \int 2\sqrt{1-\sin^2{u}}\,2\cos{u}\,du \\
	&= \int 2\cos{u}\,2\cos{u}\,du \\
	&= 4\int \cos^2{u}\,du.
\end{align*}
But this is just an antiderivative we have already found in the
exercises in chapter~11.3, namely
\begin{align*}
\int \cos^2{u}\,du &= \dfrac{u}{2} + \dfrac{1}{4} \sin 2u + C \\[3pt]
	&= \dfrac{u}{2} + \dfrac{1}{4}\cdot 2\sin u \cos u+ C \\[3pt]
	&=  \dfrac{1}{2}(u +\sin u \cos u)+ C.
\end{align*}
To find the desired antiderivative for the original function of $x$,
we now replace $u$ by its expression in terms of $x$ by inverting the
relationship: If $x=2\sin{u}$, then $\sin{u}=x/2$, and
$u=\arcsin(x/2)$.  As we found in chapter~11.1, drawing a picture
expressing the relationship between $x$ and $u$ makes it easy to
visualize the other trigonometric functions:
\begin{center}
\begin{picture}(180,80)(-10,5)
\bezier{150}(30,10)(29.2,16.89)(26.83,23.42)
\put(165,47){$x$}
\put(78,57){$2$}
\put(80,5){\makebox(0,0)[t]{$\sqrt{4-x^2}$}}
\put(32,15){$u$}
\thicklines
\put(0,10){\line(1,0){160}}
\put(160,10){\line(0,1){80}}
\put(0,10){\line(2,1){160}}
\end{picture}
\end{center}
From the picture we see that
$$
\cos{u} = \dfrac{\sqrt{4-x^2}}{2} \qquad \mbox{and} \qquad 
      \tan{u}=\dfrac{x}{\sqrt{4-x^2}}.
$$

We can now find an expression for the desired antiderivative in terms
of~$x$:
\begin{align*}
\int \sqrt{4-x^2}\,dx
   &= 2(u +\sin u \cos u)+ C \\
   &= 2\left(\arcsin\dfrac{x}{2} 
               +\dfrac{x}{2}\dfrac{\sqrt{4-x^2}}{2}\right) + C \\
   &= 2\arcsin\dfrac{x}{2}+\dfrac{x\sqrt{4-x^2}}{2}+ C.
\end{align*}
As usual, you should check this result by differentiating the
right-hand side to see that you do obtain the integrand on the left.

\medskip

Similar substitutions 
%
\mar{Some useful substitutions}
%
allow us to evaluate other integrals involving
square roots of quadratic expressions.  Here is a summary of useful
substitutions. In each case, $a$ is a positive real number.  
\medskip
\begin{center}
\begin{picture}(270,60)
	\begin{thicklines}
	\put(0,0){\framebox(270,60)
		{\begin{minipage}{200pt}
	{\hspace*{-12pt}
\begin{tabular}{lccl}
To transform  &$a^2-x^2$&  let &$x=a\,\sin{u}$; \\ [2pt]
To transform& $a^2+x^2$& let&$x=a\,\tan{u}$; \\ [2pt]
To transform& $x^2-a^2$& let &$x=a\,\sec{u}$. \\ 
\end{tabular}
}
\end{minipage}
}}
	\end{thicklines}
\end{picture}
\end{center}

\medskip
\noindent
{\bf Example 2}\quad Integrate
$$
\int \dfrac{dx}{\sqrt{x^2+9}}.
$$
If we set $x=3\tan{u}$, then $dx=3\,\sec^2{u}\,du$, and the integral becomes
\begin{align*}
\dint \dfrac{dx}{\sqrt{x^2+9}}
   &= \dint \dfrac{3\,\sec^2{u}\,du}{\sqrt{9\,\sec^2{u}}} \\
   &= \dint \sec{u}\,du \\
   &= \ln{|\sec{u} + \tan{u}|}+C,
\end{align*}
(as we saw in chapter 11.4).  To express this in terms of $x$, we
again draw a picture showing the relation between $u$ and $x$:

\begin{center}
\begin{picture}(180,80)(-10,0)
\bezier{150}(30,10)(29.2,16.89)(26.83,23.42)
\put(166,46){$x$}
\put(79,58){\makebox(0,0)[r]{$\sqrt{9 + x^2}$}}
\put(81,5){\makebox(0,0)[t]{$3$}}
\put(32,15){$u$}
\thicklines
\put(0,10){\line(1,0){160}}
\put(160,10){\line(0,1){80}}
\put(0,10){\line(2,1){160}}
%\put(20,10){\line(0,1){2.8}}
%\put(17.88,18.94){\line(1,-2){1.17}}
%\put(20,12.8){\line(-1,4){2.12}}
\end{picture}
\end{center}
From this picture we see that $\sec{u}=\sqrt{9+x^2}/3$.  Therefore
\begin{align*}
\dint \dfrac{dx}{\sqrt{x^2+9}} 
  &= \ln\left|\dfrac{\sqrt{9+x^2}}{3} +\dfrac{x}{3}\right|+C \\[3pt]
  &= \ln\left|\dfrac{\sqrt{9+x^2} +x}{3}\right|+C 
     = \ln\left|\sqrt{9+x^2} +x \right|+C',
\end{align*}
where $C'=C-\ln{3}$ is a new constant.  As usual, you should
differentiate to check that this really is the claimed antiderivative

\medskip
\noindent
{\bf Example 3}\quad Evaluate
$$
\int \dfrac{dx}{\sqrt{9x^2-16}}.
$$
We first write $\sqrt{9x^2-16}$ as
$\sqrt{9}\sqrt{x^2-(16/9)}=3\sqrt{x^2-(16/9)}$.  Using the
substitution $x=(4/3)\sec{u}$, with $dx=(4/3)\sec{u}\tan{u}\,du$ gives

\newpage

\mbox{}\vspace{-20pt}

\begin{align*}
\dint \dfrac{dx}{\sqrt{9x^2-16}}
  &= \dint \dfrac{(4/3)\sec{u}\,\tan{u}\,du}{3\sqrt{(16/9)
                     \sec^2{u}-(16/9)}} \\[6pt]
  &= \dint \dfrac{(4/3)\sec{u}\,\tan{u}\,du}{3\cdot (4/3)\tan{u}} 
                     =\dfrac{1}{3}\dint \sec{u}\,du \\[3pt]
  &= \dfrac{1}{3}\ln{| \sec{u}+\tan{u}|}+C.
\end{align*}
Again we need a picture to relate $x$ and $u$:
\begin{center}
\begin{picture}(180,90)(-10,0)
\bezier{150}(30,10)(29.2,16.89)(26.83,23.42)
\put(165,46){$\sqrt{9x^2-16}$}
\put(80,58){\makebox(0,0)[r]{$3x$}}
\put(82,5){\makebox(0,0)[t]{$4$}}
\put(32,15){$u$}
\thicklines
\put(0,10){\line(1,0){160}}
\put(160,10){\line(0,1){80}}
\put(0,10){\line(2,1){160}}
%\put(20,10){\line(0,1){2.8}}
%\put(17.88,18.94){\line(1,-2){1.17}}
%\put(20,12.8){\line(-1,4){2.12}}
\end{picture}
\end{center}
Thus $\tan{u}=\sqrt{9x^2-16}/4$, which gives
\begin{align*}
\dint \dfrac{dx}{\sqrt{9x^2-16}} 
  &= \dfrac{1}{3}\ln{| \sec{u}+\tan{u}|}+C \\
  &= \dfrac{1}{3}\ln\left|\dfrac{3x}{4} + \dfrac{\sqrt{9x^2-16}}{4}\right|+C \\
  &= \dfrac{1}{3}\ln{\left|3x +\sqrt{9x^2-16}\right|} +C',
\end{align*}
where $C'=C-(\ln{4})/3$.

As usual, you should differentiate this final expression to confirm that it really is the desired antiderivative.

\subsection{Inverse Substitution and Definite Integrals}

We saw on page~\pageref{defint} in chapter~11.2 how to use
substitution to evaluate a definite integral. When we transformed an
integral originally expressed in terms of a variable $x$ into one
expressed in terms of a variable $u$, the two integrals had the same
numerical value.  The same can be done with the inverse substitution
technique we have just been considering.  Let's see how this works.
Suppose we start with a function $f(x)$ to be integrated over an
interval $[a, b]$.  If we only knew an antiderivative $F$ for $f$, we
could easily write
$$
\int_a^b f(x)\,dx=F(b)-F(a),
$$
as usual.  In the examples we've just been considering, we found the
antiderivative for $f$ by making a substitution $x=G(u)$ for some
function $G$ and then finding an antiderivative for $f(G(u))g(u)$,
where $G'=g$.  This antiderivative we know is $F(G(u))$, where $F$ is
the function we are trying to find.  We were able to obtain $F$ by
replacing $u$ by its expression in $x$.  To do this we needed to find
the inverse function $G^{-1}$ for $G$, so that $x=G(u)$ was equivalent
to $u=G^{-1}(x)$, and $F(x)=F(G(G^{-1}(x)))$.  It is this last step we
can eliminate in calculating definite integrals.

If we want $x$ to go from $a$ to $b$, what must $u$ do? What interval
of $u$ values will get transformed to this interval of $x$ values by
$G$.  The value of $u$ such that $G(u)=a$ is $A=G^{-1}(a)$.  Similarly
the $u$ value that gets transformed to $b$ is $B=G^{-1}(b)$.  Thus
under the substitution $x=G(u)$, as $u$ goes from $A$ to $B$, $x$ will
go from $a$ to $b$.  Now look at the corresponding definite integral:
\begin{align*}
\dint_A^B f(G(u))\,g(u)\,du 
   &= \left.F(G(u))\rule{0pt}{16pt}\,\right|_A^B \\ 
   &= F(G(B))-F(G(A)) \\	
   &= F(G(G^{-1}(b)))-F(G(G^{-1}(a))) \\
   &= F(b)-F(a),
\end{align*}
which is just the desired value of the original definite integral.  To
summarize,
\begin{center}
\begin{picture}(350,100)
	\begin{thicklines}
	\put(0,0){\framebox(350,100)
		{\begin{minipage}{310pt}
	{\bf \boldmath If we make the substitution $x=G(u)$, then
$$\dint_a^b f(x)\,dx = F(b)-F(a)=\dint_A^B f(G(u))\,g(u)\,du,$$
where $A=G^{-1}(a)$ and $B=G^{-1}(b)$.
}
\end{minipage}
}}
	\end{thicklines}
\end{picture}
\end{center}
Let's look back at a couple of the preceding examples to see how this
works.

\medskip
\noindent
{\bf Example 4}\quad Evaluate
$$
\int_{-2}^2 \sqrt{4-x^2}\,dx.
$$
In Example 1 we found an antiderivative for this function by making
the substitution $x= 2\sin{u}=G(u)$.  The inverse function is then
$G^{-1}(x)=\arcsin(x/2)$. To get $x$ to go from $-2$ to 2, $u$ must go
from $G^{-1}(-2)=\arcsin(-1)=-\pi/2=A$ to
$G^{-1}(2)=\arcsin{1}=\pi/2=B$.  Then to evaluate the integral from
$x=-2$ to $x=2$, we only need to evaluate the antiderivative we found
for the $u$ integral between $u=-\pi/2$ and $u=\pi/2$.
$$
\int_{-2}^2 \sqrt{4-x^2}\,dx
  = \left.2(u+\sin{u}\,\cos{u})\rule{0pt}{16pt}\,\right|_{-\pi/2}^{\pi/2} 
  = \pi-(-\pi)=2\pi.
$$ 
(Note that this is just half the area of a circle of radius 2.  How
  could we have foreseen this result from the form of the problem?)

\medskip
\noindent
{\bf Example 5}\quad Suppose we wanted the integral
$$
\int_0^3 \dfrac{dx}{\sqrt{x^2+9}}.
$$
In Example 2 we found an antiderivative by letting $x=3\tan{u}$.  Here
$G^{-1}(x)=\arctan(x/3)$. For $x$ to go from 0 to $3$, $u$ must go
from 0 to $\arctan{1}=\pi/4$.  Using the $u$-antiderivative we found
in Example 2, we have
\begin{align*}
\dint_0^3 \dfrac{dx}{\sqrt{x^2+9}}
  &= \left.\ln|\sec{u}+\tan{u}|\rule{0pt}{16pt}\,\right|_0^{\pi/4} \\
  &= \ln{|\sqrt{2}+1|}-\ln{|1+0|}=\ln(\sqrt{2}+1).
\end{align*}


\subsection{Completing The Square}

Integrands involving terms of the form $Ax^2+Bx+C$ can always be put
in the form $A( u^2\pm b^2)$ for a suitable variable $u$ and constant
$b$.  The technique for doing this is the standard method of {\bf
  completing the square}\index{completing the square}:
\begin{align*}
Ax^2+Bx+C
  &= A\left(x^2 + \dfrac{B}{A}x\right)+C \\
  &= A\left(x^2 + \dfrac{B}{A}x + \dfrac{B^2}{4A^2}\right)+C 
         - \dfrac{B^2}{4A} \\
  &= A\left(x + \dfrac{B}{2A}\right)^{\!2} + \dfrac{4AC-B^2}{4A}.
\end{align*}
The substitutions
$$
u = x+ \dfrac{B}{2A} \qquad \mbox{and} \qquad 
  b = \dfrac{\sqrt{|4AC-B^2|}}{2A}
$$
then transform the problem to a form where we can use the techniques
already developed.  The following examples should make this clear.

\medskip
\noindent
{\bf Example 6}\quad Consider the integral
$$
\int \dfrac{dx}{x^2+4x+5}.
$$
This may not immediately remind us of anything we've seen before.  But
if we rewrite it in the form
$$
\int \dfrac{dx}{(x^2+4x+4)+1}=\int \dfrac{dx}{(x+2)^2+1},
$$ 
it now begins to resemble something involving an arctangent.  In fact,
if we make the substitution $u=x+2$, so $du=dx$, we can write
\begin{align*}
\dint \dfrac{dx}{x^2+4x+5} 
        &= \dint \dfrac{du}{u^2+1} \\
	&= \arctan{u}+C \\
	&= \arctan(x+2) +C.
\end{align*}

\medskip
\noindent
{\bf Example 7}\quad The technique of completing the square even works
for expressions we could have factored directly, if we had noticed:
\begin{align*}
\dint \dfrac{dx}{x^2+4x+3} 
   &= \dint \dfrac{dx}{(x+2)^2-1} \\[3pt]
   &= \dint \dfrac{dx}{(x+2 -1)(x+2+1)} \\[3pt]
   &= \dint \dfrac{dx}{(x+1)(x+3)} \\[3pt]
   &= \dfrac{1}{2}\dint \dfrac{dx}{x+1}
           - \dfrac{1}{2}\dint \dfrac{dx}{x+3} \\[5pt]
   &= \dfrac{1}{2}\ln\left|\dfrac{x+1}{x+3}\right| +C.
\end{align*}


\medskip
\noindent
{\bf Example 8}\quad Evaluate
$$\int \dfrac{dx}{\sqrt{6x-x^2}}. $$
Note that $6x-x^2=-(x^2-6x)=-(x-3)^2+9=9-(x-3)^2$.  If we now substitute $x-3 =3u$, with $dx=3du$, we get
\begin{align*}
\dint \dfrac{dx}{\sqrt{6x-x^2}}
        &= \dint \dfrac{dx}{\sqrt{9-(x-3)^2}}\\
	&= \dint \dfrac{3 \, du}{\sqrt{9-9u^2}}
              = \dint \dfrac{du}{\sqrt{1-u^2}} \\
	&= \arcsin{u}+C \\
	&= \arcsin \dfrac{x-3}{3} +C.
\end{align*}

\subsection{Trigonometric Polynomials}

A {\bf trigonometric polynomial}\index{polynomial!trigonometric} is
any sum of constant multiples of products of trigonometric functions.
The preceding techniques have shown some cases where such
trigonometric polynomials can arise, even though the original problem
had no apparent reference to trigonometric functions.  There are many
different ways of breaking trigonometric polynomials down into special
cases which can then be integrated.  We will develop one way which has
the virtue of using few special cases, so that it can be used fairly
automatically.  It also introduces a powerful tool---that of {\bf
  reduction formula}---which can be used to generate mathematical
results interesting in their own right.  One example is the striking
representation of $\pi$ derived in chapter~12.1.  Other examples are
developed in the exercises at the end of this section.

Since every trigonometric function is expressible in terms of sines
and cosines, any trigonometric polynomial can be written as a sum of
terms of the form $c\,\sin^m{x}\,\cos^n{x}$ where $c$ is a constant
and $m$ and $n$ are integers---positive, negative, or 0.  For
instance, $5\,\sec^2{x}\,\tan^5{x}$ can be rewritten as
$5\,\sin^5{x}\,\cos^{-7}{x}$. To find antiderivatives for
trigonometric polynomials, it therefore suffices to be able to
evaluate integrals of the form
$$
\int \sin^m{x}\cos^n{x}\,dx.
$$

We will see how to find antiderivatives for functions of this sort by
breaking the problem into a series of special cases:

\medskip
\begin{tabular}{rl}
Category I&either $m\geq 0$ or $n\geq 0$ (or both) \\
Case 1&$m=1$ or $n=1$ \\
Case 2&$m=0$ or $n=0$ \\
Category II& $m$ and $n$ both negative \\
\end{tabular}


\subsubsection{Category I: Either \boldmath $m\geq 0$ or $n\geq 0$ (or both)}

Assume for the sake of explicitness that $m\geq 0$. We can then use
the identity $\sin^2{x}=1-\cos^2{x}$ to replace $\sin^m{x}$ entirely
by cosine terms if $m$ is even, or to replace all but one of the sine
terms by cosines if $m$ is odd.  A similar replacement can be made if
$n\geq 0$.


\medskip
\noindent
{\bf Example 9}\quad
\begin{align*}
\sin^4{x}\,\cos^6{x}
  &= (1-\cos^2{x})^2\cdot\cos^6{x} \\
  &= (1-2\,\cos^2{x}+\cos^4{x})\cdot\cos^6{x} \\
  &= \cos^6{x}-2\,\cos^8{x}+\cos^{10}{x}.
\end{align*}
(Note that in this example we could just as well have expressed
$\cos^6{x}$ entirely in terms of $\sin{x}$.)

\medskip
\noindent
{\bf Example 10}\quad
\begin{align*}
\sin^3{x}\,\cos^{-8}{x}
  &= \sin{x}\cdot (1-\cos^2{x})\cdot\cos^{-8}{x} \\
  &= \sin{x}\,\cos^{-8}{x}-\sin{x}\,\cos^{-6}{x}.
\end{align*}

\medskip
\noindent
{\bf Example 11}\quad
\begin{align*}
\sin^{-7}{x}\,\cos^{7}{x}
  &= \sin^{-7}{x}\cdot (1-\sin^2{x})^3\cdot\cos{x} \\
  &= \sin^{-7}{x}\,\cos{x}-3\,\sin^{-5}{x}\,\cos{x}+3\,\sin^{-3}{x}\,\cos{x} \\
  & \mbox{}\quad -\sin^{-1}{x}\,\cos{x}.
\end{align*}

\medskip
We can thus reduce any problem in Category I to one of two special cases:
\begin{center}
\begin{tabular}{rl}
Case 1&$m=1$ or $n=1$ \\
Case 2&$m=0$ or $n=0$ \\
\end{tabular}
\end{center}
We will now see how to find antiderivatives for these cases.

\medskip
\noindent
{\bf Case 1: \boldmath $m=1$ or $n=1$}\quad Since the two
possibilities are analogous, we will consider the case with
$n=1$. Then $m$ can be any real number at all, not necessarily an
integer.  We make the substitution $u=\sin{x}$, so that
$du=\cos{x}\,dx$, and
\[
\int \sin^m{x}\cos{x}\,dx = \int u^m\,du = 
\begin{cases}
\dfrac{1}{m+1} \,u^{m+1} + C & \mbox{if $m\neq -1$}, \\[10pt]
\ln{|u|} + C & \mbox{if $m=-1$}.
\end{cases}
\]
Replacing $u$ by its expression in $x$ we have the antiderivative:

\begin{center}
\begin{picture}(340,70)
	\begin{thicklines}
	\put(0,0){\framebox(340,70)
{ $\dint \sin^m{x}\cos{x}\,dx=\begin{cases}
\dfrac{1}{m+1} \sin^{m+1}{x} + C & \mbox{if $m\neq -1$}, \\[10pt]
\ln{|\sin{x}|} + C & \mbox{if $m=-1$}.
\end{cases}
$}}
	\end{thicklines}
\end{picture}
\end{center}

\noindent
{\bf Remark:}  The instance $m=-1$ in this 
%
\mar{The antiderivative\\ of the cotangent}\index{cotangent!antiderivative of}
%
case is worth singling out, as it gives us an antiderivative for
$\cot{x}$:
$$
\int \cot{x}\,dx=\int \dfrac{\cos{x}}{\sin{x}}\,dx =  \ln|\sin{x}|+C.
$$

Integrals where $m=1$ are handled in a completely analogous fashion.
You should check that
\begin{center}
\begin{picture}(340,70)
	\begin{thicklines}
	\put(0,0){\framebox(340,70)
{ $\dint \cos^n{x}\sin{x}\,dx=\begin{cases}
\dfrac{-1}{n+1} \cos^{n+1}{x} + C & \mbox{if $n\neq -1$}, \\[10pt]
-\ln{|\cos{x}|} + C & \mbox{if $n=-1$}.
\end{cases}
$}}
	\end{thicklines}
\end{picture}
\end{center}

\noindent
{\bf Remark:} Notice that  $n=-1$ gives us 
%
\mar{The antiderivative\\ of the tangent}\index{tangent!antiderivative of}
%
an antiderivative for $\tan{x}$:
$$ 
\int \tan{x}\,dx=\int \dfrac{\sin{x}}{\cos{x}}\,dx = -\ln|\cos{x}|+C.
$$
\label{tander}

\noindent
{\bf Case II: \boldmath $m=0$ or $n=0$}\quad Again the two possibilities are
analogous, so we will look at instances where $n=0$.  There are a
number of clever ways for dealing with antiderivatives of functions of
this form, many of them depending on special subcases according to
whether $m$ is even or odd, positive or negative, etc.  We will
develop a single method which deals with all cases in the same way.

Think of $\sin^n{x}$ as $\sin^{n-1}{x}\cdot\sin{x}$ and use
integration by parts with
$$
F(x)=\sin^{n-1}{x}\qquad\mbox{ and }\qquad g(x)=\sin{x};
$$
then
$$
f(x)=(n-1)\sin^{n-2}{x}\,\cos{x}\qquad\mbox{ and }\qquad G(x)=-\cos{x}.
$$
Therefore
$$
\int \sin^n{x}\,dx
    =-\sin^{n-1}{x}\,\cos{x}+(n-1)\int \sin^{n-2}{x}\,\cos^2{x}\,dx.
$$ 

Now since $\cos^2{x}=1-\sin^2{x}$, we can rewrite the integral on the
right-hand side as
$$
\int \sin^{n-2}{x}\,\cos^2{x}\,dx.=\int \sin^{n-2}{x}\,dx-\int \sin^n{x}\,dx
$$
---an expression involving the original integral we are trying to
evaluate!  If we now substitute this expression in our original
equation and bring all the terms involving $ \sin^n{x}$ over to the
left-hand side, we have
$$
n\int  \sin^n{x}\,dx 
   = -\sin^{n-1}{x}\,\cos{x}+(n-1)\int \sin^{n-2}{x}\,dx,
$$
so that
\begin{center}
\begin{picture}(306,50)\label{reduction}
	\begin{thicklines}
	\put(0,0){\framebox(306,50)
{ $\dint \sin^n{x}\,dx=\dfrac{-1}{n}\sin^{n-1}{x}\,\cos{x}
     + \dfrac{n-1}{n}\dint \sin^{n-2}{x}\,dx.
$}}
	\end{thicklines}
\end{picture}
\end{center}

We thus have 
%
\mar{A reduction formula}
%
a {\bf reduction formula}\index{reduction formula!for integrating
  $\sin^n{x}$} which reduces the problem of finding an antiderivative
for $\sin^n{x}$ to the problem of finding an antiderivative for
$\sin^{n-2}{x}$.  This in turn can be reduced to finding an
antiderivative for $\sin^{n-4}{x}$, and so on, until we get down to
having to find an antiderivative for $\sin{x}$ (if $n$ is odd), or for
1 (if $n$ is even).

\medskip
\noindent
{\bf Example 12}
\begin{align*}
\dint \sin^5{x}\,dx
  &= \dfrac{-1}{5}\sin^4{x}\cos{x}+\dfrac{4}{5}\dint \sin^3{x}\,dx \\
  &= \dfrac{-1}{5}\sin^4{x}\cos{x}
       +\dfrac{4}{5}\left(\dfrac{-1}{3}\sin^2{x}\cos{x}
       +\dfrac{2}{3}\dint \sin{x}\,dx \,\right) \\
  &= \dfrac{-1}{5}\sin^4{x}\cos{x}-\dfrac{4}{15}\sin^2{x}\cos{x}
       -\dfrac{8}{15} \cos{x}+C.
\end{align*}
Check this answer by taking the derivative of the right-hand side.  To
show that this derivative really is equal to the integrand on the
left, you will need to express all the cosines in terms of sines.

\medskip
\noindent
{\bf Example 13}\quad If we let $n=2$, we quickly get the
antiderivative for $\sin^2{x}$ that we've needed at several points
already:
\begin{align*}
\dint \sin^2{x}\,dx 
  &= \dfrac{-1}{2}\sin{x}\cos{x}+\dfrac{1}{2}\dint 1\,dx \\
  &= \dfrac{x}{2}- \dfrac{1}{2}\sin{x}\cos{x}.
\end{align*}

\medskip
The reduction formula as stated is 
%
\mar{In its current form,\\ the reduction formula works best for $n>0$}
%
most convenient for $n>0$, although it is true for any number $n\neq
0$.  For if $n<0$, though, we want to {\em increase} the exponent,
replacing a problem of finding an antiderivative for $\sin^n{x}$ by a
problem where the exponent is less negative.  We can do this by
rearranging the formula as
$$
\int \sin^{n-2}{x}\,dx
  = \dfrac{n}{n-1}\int \sin^n{x}\,dx+ \dfrac{1}{n-1}\sin^{n-1}{x}\cos{x}.
$$ 
Since we are interested in negative exponents, 
%
\mar{The reduction formula for negative exponents}
%
call $n-2$ by a new name, $-k$.  But if $n-2=-k$, then $n=-k+2$, and
we can rewrite our formula as
\begin{center}
\begin{picture}(360,50)
	\begin{thicklines}
	\put(0,0){\framebox(360,50)
{ $\dint \sin^{-k}{x}\,dx=-\dfrac{1}{k-1}\sin^{-(k-1)}{x}\cos{x}+\dfrac{k-2}{k-1}\int \sin^{-(k-2)}{x}\,dx.
$}}
	\end{thicklines}
\end{picture}
\end{center}

With this formula we can reduce the problem of finding an
antiderivative for $\sin^{-k}{x}$ to the problem of finding an
antiderivative for $\sin^{-k+2}{x}$.  This in turn can be reduced to
finding an antiderivative for $\sin^{-k+4}{x}$, and so on, until we
get up to having to find an antiderivative for $\sin^{-1}{x}$ (if $k$
is odd), or for 1 (if $k$ is even).  All we need, then, is an
antiderivative for $\sin^{-1}{x}$.  But $\sin^{-1}{x}=\csc{x}$, and in
chapter~11.4 (page~\pageref{trick}) we found that
$$
\int \csc{x}\,dx=\int \sin^{-1}{x}\,dx=-\ln{|\csc{x}+\cot{x}|}+C.
$$

We can now handle antiderivatives for any negative integer exponent of
the sine function.

\medskip
\noindent
{\bf Example 14} \quad We can check this formula by trying $k=2$, which will give us the antiderivative of $\csc^2{x}$:
\begin{align*}
\dint \csc^2{x}
  &= \dint \sin^{-2}{x}\,dx \\
  &= -\dfrac{1}{1}\sin^{-1}{x}\,\cos{x}+\dfrac{0}{1}\int \sin^0{x}\,dx \\
  &= -\sin^{-1}{x}\,\cos{x}+C=-\cot{x}+C,
\end{align*}
as it should.

\medskip
\noindent
{\bf Example 15} \quad
\begin{align*}
\dint \sin^{-3}{x}\,dx 
  &= -\dfrac{1}{2}\sin^{-2}{x}\,\cos{x}+\dfrac{1}{2}\int \sin^{-1}{x}\,dx \\
  &= \dfrac{1}{2}\left(-\sin^{-2}{x}\cos{x}-\ln{|\csc{x}+\cot{x}|}\right)+C.
\end{align*}

In the exercises you are asked to derive the following reduction
formulas for the cosine function\index{reduction formula!for cosine}:
\begin{center}
\begin{picture}(360,90)\label{cosred}
	\begin{thicklines}
	\put(0,0){\framebox(360,90)
		{\begin{minipage}{330pt}
\mbox{}\vspace{-15pt}
\begin{gather*}
\int \cos^m{x}\,dx = \dfrac{1}{m}\cos^{m-1}{x}\,\sin{x} 
          + \dfrac{m-1}{m}\int \cos^{m-2}{x}\,dx, \\[12pt]
\int \cos^{-m}{x}\,dx = \dfrac{1}{m-1}\cos^{-m+1}{x}\,\sin{x} 
          + \dfrac{m-2}{m-1}\int \cos^{-m+2}{x}\,dx.
\end{gather*}
\end{minipage}
}}
	\end{thicklines}
\end{picture}
\end{center}


\subsubsection{Category II: Both \boldmath $m<0$ and $n<0$}

If we divide the identity $\cos^2{x}+\sin^2{x}=1$ by
$\sin^2{x}\,\cos^2{x}$, we get the identity
$$
\sin^{-2}{x}+\cos^{-2}{x}=\sin^{-2}{x}\,\cos^{-2}{x}.
$$ 
We will now use this identity to express anything of the form
$\cos^{-r}{x}\,\sin^{-s}{x}$ (where $r>0$ and $s>0$) as a sum of terms
of the form $\sin^{-h}{x}$, or $\cos^{-i}{x}$, or
$\sin{x}\,\cos^{-j}{x}$, or $\sin^{-k}{x}\,\cos{x}$.  Since we learned
how to find antiderivatives for expressions like these in the previous
cases, we will then be done.

The trick in transforming $cos^{-r}{x}\,\sin^{-s}{x}$ to the desired
form is to multiply by $(\cos{x}\,\cos^{-1}{x})$ or
$(\sin{x}\,\sin^{-1}{x})$ as needed so that both the sine and the
cosine terms appear to {\em even} negative exponents.  Then simply
keep using the identity above until there's nothing left to use it
on. The following three examples should make clear how the reduction
then works.

\medskip
\noindent
{\bf Example 16} \quad($r$ and $s$ both even already)
\begin{align*}
\sin^{-4}{x}\,\cos^{-6}{x}
  &= \left( \sin^{-2}{x}\,\cos^{-2}{x}\right)^2\,\cos^{-2}{x} \\
  &= \left( \sin^{-2}{x}+\cos^{-2}{x}\right)^2\,\cos^{-2}{x} \\
  &= (\sin^{-4}{x}+2\sin^{-2}{x}\,\cos^{-2}{x}+\cos^{-4}{x})\,\cos^{-2}{x} \\
  &= (\sin^{-4}{x}+2(\sin^{-2}{x}+\cos^{-2}{x})+\cos^{-4}{x})\,\cos^{-2}{x} \\
  &= \sin^{-4}{x}\,\cos^{-2}{x}+2\sin^{-2}{x}\,\cos^{-2}{x}+2\cos^{-4}{x} \\
  & \mbox{}\qquad +\cos^{-6}{x} \\
  &= \sin^{-2}{x}\,(\sin^{-2}{x}+\cos^{-2}{x})+2(\sin^{-2}{x}+\cos^{-2}{x}) \\
  & \mbox{}\qquad +2\cos^{-4}{x}+\cos^{-6}{x} \\
  &= \sin^{-4}{x}+\sin^{-2}{x}\,\cos^{-2}{x}+2\sin^{-2}{x}+2\cos^{-2}{x} \\
  & \mbox{}\qquad +2\cos^{-4}{x}+\cos^{-6}{x} \\
  &= \sin^{-4}{x}+( \sin^{-2}{x}+\cos^{-2}{x}) +2\sin^{-2}{x} \\
  & \mbox{}\qquad + 2\cos^{-2}{x}+ 2\cos^{-4}{x}+ \cos^{-6}{x} \\
  &= \sin^{-4}{x}+3\sin^{-2}{x}+ 3\cos^{-2}{x}+2\cos^{-4}{x}+ \cos^{-6}{x}.
\end{align*}
While this process is tedious, it requires little thought---you simply
replace $\sin^{-2}{x}\,\cos^{-2}{x}$ with $\sin^{-2}{x}+\cos^{-2}{x}$
at every opportunity until there is no negative-exponent sine term
multiplying any negative-exponent cosine term.  We will use this
result to demonstrate how to deal with cases where either $r$ or $s$
(or both) is odd.

\medskip

\noindent
{\bf Example 17} \quad($r$ even and $s$ odd)\label{ex17}
\begin{align*}
\sin^{-4}{x}\,\cos^{-5}{x}
  &= \sin^{-4}{x}\,\cos^{-6}{x}\,\cos{x} \\
  &= \sin^{-4}{x}\,\cos{x}+3\sin^{-2}{x}\,\cos{x}+ 3\cos^{-1}{x} \\
  &  \mbox{}\qquad +2\cos^{-3}{x}+ \cos^{-5}{x}
\end{align*}

\medskip
\noindent
{\bf Example 18} \quad(both $r$ and $s$ odd)
\begin{align*}
\sin^{-3}{x}\,\cos^{-5}{x}
  &= \sin{x}\,\sin^{-4}{x}\,\cos^{-6}{x}\,\cos{x} \\
  &= \sin^{-3}{x}\,\cos{x}+3\sin^{-1}{x}\,\cos{x}+ 3\sin{x}\,\cos^{-1}{x} \\
  & \mbox{}\qquad +2\sin{x}\,\cos^{-3}{x}+ \sin{x}\,\cos^{-5}{x}
\end{align*}


\subsection{Exercises}

\vspace{-10pt}

\num Find the following antiderivatives ($a$ is a positive constant):

\noindent
\begin{minipage}[t]{100pt}
\sub $\dint \dfrac{dx}{\sqrt{1-4x^2}}$

\sub $\dint \dfrac{dx}{\sqrt{1-4x^2}}$

\sub $\dint \dfrac{dx}{\sqrt{4+x^2}}$

\sub $\dint \dfrac{x\,dx}{\sqrt{4+x^2}}$
\end{minipage}
\hfill
\begin{minipage}[t]{100pt}
\sub $\dint \dfrac{dx}{(a^2-x^2)^{3/2}}$

\sub $\dint \dfrac{dx}{4+x^2}$

\sub $\dint \dfrac{x\,dx}{4+x^2}$
\end{minipage}
\hfill
\begin{minipage}[t]{100pt}
\sub $\dint \dfrac{dx}{x\sqrt{4+x^2}}$

\sub $\dint \dfrac{x\,dx}{\sqrt{x^2-a^2}}$

\sub $\dint \dfrac{dx}{(a^2+x^2)^2}$
\end{minipage}
\hfill
\mbox{}

\num Evaluate the following integrals:

\noindent
\begin{minipage}[t]{150pt}
\sub $\dint_1^{-1} \dfrac{dx}{4-x^2}$
\vspace*{4pt}
\sub $\dint_1^2 \sqrt{x^2-1}\,dx$
\vspace*{4pt}
\sub $\dint_0^{\pi/3} x\,\sec^2{x}\,dx$
\end{minipage}
\hfill
\begin{minipage}[t]{2.4in}
\sub $\dint_0^1 \dfrac{dx}{(2-x^2)^{3/2}}$
\vspace*{4pt}
\sub $\dint_0^\infty \dfrac{dx}{9+x^2}$
\vspace*{4pt}
\sub $\dint_a^{2a} x^3\sqrt{x^2-a^2}\,dx$
\end{minipage}
\hfill
\mbox{}

\num Sketch the ellipse $\dfrac{x^2}{a^2}+ \dfrac{y^2}{b^2} = 1$,
labelling the coordinates of the points where it crosses the $x$-axis
and the $y$-axis.  Prove that the area of this ellipse is $\pi ab$.

\num Find the following antiderivatives:

\noindent
\begin{minipage}[t]{150pt}
\sub $\dint \dfrac{dx}{\sqrt{x^2-2x-8}}$
\vspace*{4pt}
\sub $\dint \dfrac{dx}{x^2+6x+10}$
\vspace*{4pt}
\sub $\dint \dfrac{dx}{\sqrt{x^2+6x+8}}$
\vspace*{4pt}
\sub $\dint \dfrac{x\,dx}{x^2+4x+5}$
\end{minipage} 
\hfill
\begin{minipage}[t]{150pt}
\sub $\dint \dfrac{x\,dx}{\sqrt{5+4x-x^2}}$
\vspace*{4pt}
\sub $\dint \dfrac{(2x+7)\,dx}{4x^2+4x+5}$
\vspace*{4pt}
\sub $\dint \dfrac{(4x-3)\,dx}{\sqrt{-x^2-2x}}$
\vspace*{4pt}
\sub $\dint \dfrac{dx}{(a^2-x^2-2x)^2}$
\end{minipage}
\hfill
\mbox{}

\numa If $x= a\,\sec{u}$, where $a$ is a constant, draw a right
triangle containing an angle $u$ with lengths of sides specified to
reflect this relation between $x$, $a$, and $u$.

\sub In terms of $x$ and~$a$, what is $\sin{u}$?  What is $\cos{u}$?
What is $\tan{u}$?

\num Evaluate the following:

\noindent
\begin{minipage}[t]{150pt}
\sub $\dint \dfrac{dx}{\sin{x}\,\cos{x}}$
\vspace*{4pt}
\sub $\dint \cos^3{x}\,\sin^{-4}{x}\,dx$
\vspace*{4pt}
\sub $\dint \csc^4{x}\,\cot^2{x}\,dx$
\vspace*{4pt}
\sub $\dint \sin{3x}\,\cot{3x}\,dx$
\vspace*{4pt}
\sub $\dint_0^{\pi/2} \sin^n{x}\,\cos^3{x}\,dx$
\end{minipage} 
\hfill
\begin{minipage}[t]{150pt}
\sub $\dint \tan^5{x}\,dx$
\vspace*{4pt}
\sub $\dint \dfrac{\sin^3{5x}\,dx}{\sqrt[3]{\cos{5x}}}$
\vspace*{4pt}
\sub $\dint \dfrac{\cos^3(\ln{x})\,dx}{x}$
\vspace*{4pt}
\sub $\dint \sec^4{x}\,\ln(\tan{x})\,dx$
\vspace*{4pt}
\sub $\dint_0^{a/2} \dfrac{dx}{(a^2-x^2)^{3/2}}$
\end{minipage}
\hfill
\mbox{}

\num Use the analysis of Example 17 (page~\pageref{ex17}) to find an
antiderivative for $\sin^{-4}{x}\,\cos^{-5}{x}$.

\subsubsection{Reduction formulas}
\vspace{-10pt}

\num Derive the reduction formulas for the cosine function given on
page~\pageref{cosred}.

\numa By writing $\tan^n{x}=\tan^{n-2}{x}(\sec^2{x}-1)$, get a
reduction formula which expresses $\dint \tan^n{x}\,dx$ in terms of
$\dint \tan^{n-2}{x}\,dx$.  

\sub Use this evaluation formula to find $\dint \tan^6{x}\,dx.$

\sub Show that
\[
\int_0^{\pi/4}\tan^n{x}\,dx =
\begin{cases}
\dfrac{1}{n-1}- \dfrac{1}{n-3}+\cdots\pm \dfrac{1}{3}\mp 1\pm {\pi/4} 
     &\text{if $n$ is even}, \\[10pt]
\dfrac{1}{n-1}- \dfrac{1}{n-3}+\cdots \pm \dfrac{1}{4}\mp \dfrac{1}{2}
             \pm \dfrac{1}{2}\ln{2} 
     &\text{if $n$ is odd}.
\end{cases}
\]

\sub Give a clear argument why \quad $\displaystyle
\lim_{n\rightarrow\infty}\int_0^{\pi/4} \tan^n{x}\,dx=0$.

\sub Prove that
\[
\lim_{k\rightarrow\infty}\left(1-\dfrac{1}{3}+\dfrac{1}{5}
      -\cdots\pm \dfrac{1}{2k+1}\right)=\dfrac{\pi}{4},
\]
and
\[
\lim_{k\rightarrow\infty}\left(1-\dfrac{1}{2}+\dfrac{1}{3}
       -\cdots\pm \dfrac{1}{k}\right)=\ln{2}.
\]

\numa By writing $\sec^n{x}$ as $\sec^{n-2}{x}\,\sec^2{x}$ and using
integration by parts, get a reduction formula which expresses $\dint
\sec^n{x}\,dx$ in terms of $\dint \sec^{n-2}{x}\,dx$.

\sub Since $\sec{x}=\cos^{-1}{x}$, the formula you got in part (a)
could also have been obtained from the reduction formula for cosines
on page~\pageref{cosred}.  Try it and see if the formulas are in fact
the same.

\numa Find a reduction formula that expresses
\[
\dint x^n\,e^x\,dx \ \quad \text{in terms of} \quad \dint x^{n-1}\,e^x\,dx.
\]

\sub Using the results of part (a), show that
$$
\dfrac{1}{n!}\int_0^t x^n\,e^x\,dx 
   = e^t\left(\dfrac{t^n}{n!} - \dfrac{t^{n-1}}{(n-1)!} 
       +\dfrac{t^{n-2}}{(n-2)!} - \cdots \pm \dfrac{t^2}{2!} 
       \mp t\pm 1\right)\mp 1.
$$

\sub Explain why, for a fixed value of $t$,
$$\lim_{n\rightarrow\infty} \dfrac{1}{n!}\int_0^t x^n\,e^x\,dx =0.$$

\sub Prove that
$$
\lim_{n\rightarrow\infty}\left(1-t+\dfrac{t^2}{2!}
    -\dfrac{t^3}{3!}+\dfrac{t^4}{4!}
    -\cdots \pm \dfrac{t^n}{n!}\right)=e^{-t}.
$$

\numa Find a reduction formula expressing 
\[
\dint \dfrac{dx}{(1+x^2)^n} \quad \text{in terms of} \quad
   \dint \dfrac{dx}{(1+x^2)^{n-1}};
\]
you can do this using integration by parts, or you can use a
trigonometric substitution.\label{1plusx2}

\sub What is the exact value of $\dint_0^1 \dfrac{dx}{(1+x^2)^5}$?

\num Our approach to integrating trigonometric polynomials was to
express everything in the form $\sin^m{x}\,\cos^n{x}$.  We can just as
readily express everything in the form $\sec^j{x}\tan^k{x}$ ($j$ and
$k$ integers---positive, negative, or 0), and develop our technique by
dealing with various cases of this.  See if you can work out the
details, trying to parallel the approach developed in the text using
sines and cosines as our basic functions.

%%%%%%%%%%%%%%%%%%%%%%%%%%%%%%%%%%%%%%%%%%%%%%%%%%%%%%%%%%%%%%%%%%%%%%%%
%                                                                      %
%                               Section 6                              %
%                                                                      %
%%%%%%%%%%%%%%%%%%%%%%%%%%%%%%%%%%%%%%%%%%%%%%%%%%%%%%%%%%%%%%%%%%%%%%%%

\section{Simpson's Rule}

This chapter has concentrated on formulas for antiderivatives, 
%
\mar{A return to\\ numerical methods}	
%
because a formula conveys compactly a lot of information.  However,
you must not lose sight of the fact that most antiderivatives cannot
be found by such analytic methods.  The integrand may be a data
function, for instance, and thus have no formula.  And even when the
integrand is given by a formula, there may be no formula for the
antiderivative itself.  One possibility in such cases is to {\bf
  approximate} such a function by a function---such as a
polynomial---for which we can readily find an antiderivative.  In
chapter 10 we saw some methods for doing this.  In chapter~12 we
introduce Fourier series, providing another family of approximating
functions for which antiderivatives can be readily obtained. Another
approach is to find a desired definite integral using approximating
rectangles, as we did in chapter 6.
	
In any case, numerical methods are inescapable, but accurate results
require many calculations.  This takes time---even on a modern
high-speed computer.  A numerical method is said to be {\bf
  efficient}\index{efficient computation} if it gets accurate results
quickly, that is, with relatively few calculations.  In chapter~6 we
saw that {\em midpoint\/} Riemann sums are much more efficient than
left or right {\em endpoint\/} Riemann sums.
%
\mar{Efficient numerical integration}	
%
We will look at these and other methods in detail in this section.
The most efficient method we will develop is called Simpson's rule.

\newpage

\subsection{The Trapezoid Rule}
\index{trapezoid rule}

We interpret the integral
	$$
	\int_a^b f(x) \, dx
	$$
as the area under the graph $y = f(x)$ between $x = a$ and $x = b$.
We interpret a Riemann sum as the total area of a collection of
rectangles that approximate the area under the graph.
%
\mar{Replace rectangles\\ by trapezoids}
%
The tops of the rectangles are level, and they represent the graph of
step function.  Clearly, we get a better approximation to the graph by
using slanted lines.  They form the tops of a sequence of {\bf
  trapezoids}\index{trapezoid} that approximate the area under the
graph.

\vspace{170pt}

\special{psfile=../ps/calc2/trap1.ps}

Let's figure out the areas of these trapezoids.  
%
\mar{Each trapezoid is sandwiched between two rectangles\ldots}	
%
They are related in a simple way to the rectangles that we would
construct at the right and left endpoints to calculate Riemann sums.
To see the relation, let's take a closer look at a single single
trapezoid.  

\vspace{130pt}

\special{psfile=../ps/calc2/trap2.ps} 

\newpage

\noindent 
It is sandwiched between  two rectangles, one taller and one shorter.  
%
\mar{\ldots whose average area equals the area of the trapezoid}
%
In our picture the height of the taller rectangle is $f(\mbox{right
  endpoint})$.  We will call it the {\em right rectangle\/}. The
height of the shorter rectangle is $f(\mbox{left endpoint})$.  We will
call it the {\em left rectangle\/}.  For other trapezoids the left
rectangle may be the taller one.  In any case, the trapezoid is
exactly half-way between the two rectangles in size, and thus its area
is the {\em average\/} of the areas of the rectangles:
	$$
	\mbox{area trapezoid} = \tfrac{1}{2} \left( 
		\mbox{area left rectangle} + \mbox{area right rectangle}
		\right).
	$$
If we sum over the areas of all of the trapezoids, 
%
\mar{The trapezoidal approximation is the average of left and right Riemann sums} 
%
the areas of the right rectangles sum to the right Riemann sum, and
similarly for the left rectangles.  In follows that\index{approximation!trapezoidal} \index{left Riemann sum} \index{right Riemann sum}
\[
\mbox{trapezoidal approximation}
  = \tfrac{1}{2}( \mbox{left Riemann sum} + 
			\mbox{right Riemann sum}).
\]

The pictures make it clear that the trapezoidal approximation should
be significantly better than either a left or a right Riemann sum.  To
test this numerically, let's get numerical estimates for the integral
	$$
\int_1^3 \dfrac{1}{x}\,dx = \ln{3} = 1.098612288668\ldots .
	$$
(The relation between the trapezoid approximation and the left and
right Riemann sums holds for any choice $\Delta x_k$ of subintervals.
\mar{Comparing approximations} However, we will use equal subintervals
to make the calculations simpler.)  Here is how our four main
estimates compare when we use 100 subintervals.  \label{trap}
\begin{center}
\begin{tabular}{ccr}
	$n = 100$ & approximation & \multicolumn{1}{c}{error} \\ [2pt]
	\hline \\ [-9pt]
	right & 1.09197525 & $6.63\times10^{-3}$\\
	left & 1.10530858 & $-6.69\times10^{-3}$\\
	midpoint & 1.09859747 & $1.48\times10^{-5}$\\
	trapezoidal & 1.09864191 & $-2.90\times10^{-5}$
\end{tabular}
\end{center}
The figures in this table are calculated to 8 decimal places, and the
column marked {\em error\/} is the difference
	$$
1.09861229 - \mbox{approximation}, 
	$$
so that the error is negative if the approximation is too large.  The
left Riemann sum is too large, for instance.

Before we comment on the differences between the estimates, let's
gather more data.  Here are the calculations for 1000 subintervals.
Note that the midpoint and trapezoidal approximations are more than
1000 times better than the left or right Riemann sums!
\begin{center}
\begin{tabular}{ccr}
	$n = 1000$ & approximation & \multicolumn{1}{c}{error} \\ [2pt]
	\hline \\ [-9pt]
	right & 1.09794591 & $6.663\times10^{-4}$\\
	left & 1.09927925 & $-6.669\times10^{-4}$\\
	midpoint & 1.09861214 & $1.4\times10^{-7}$\\
	trapezoidal & 1.09861258 & $-2.9\times10^{-7}$
\end{tabular}
\end{center}

We expected the trapezoidal approximation   
%
\mar{A surprise:\\ the midpoint approximation is even better than the trapezoidal}	
%
to be better than either the right or left Riemann sum.  The
surprising observation is that the midpoint Riemann sum is even
better!  In fact, it appears that the midpoint Riemann sum has only
{\em half\/} the error of the trapezoidal approximation.

The figures below explain geometrically why the midpoint approximation
is better than the trapezoidal.  The first step is shown on the left.
Take a midpoint rectangle (whose height is $f(\mbox{midpoint})$), and
rotate the top edge around the midpoint until it is tangent to the
graph of $y = f(x)$.  Call this a {\bf midpoint
  trapezoid}.\index{midpoint trapezoid} Notice that the trapezoid has
the same area as the rectangle.


\mbox{}
\vspace{100pt}

\special{psfile=../ps/calc2/trap3.ps}

\vspace{20pt}
	
The second step is to compare 
%
\mar{Shading depicts\\ the errors}	
%
the midpoint trapezoid to the one used in the trapezoidal
approximation.  This is done on the right.  The error coming from the
midpoint trapezoid is shaded light gray, while the error from the
trapezoidal approximation is dark gray.  The midpoint trapezoid is the
better approximation.  Since the midpoint rectangle has the same area
as the midpoint trapezoid, we now see why the midpoint Riemann sum is
more accurate than the trapezoidal approximation.  This picture also
explains why the errors of the two approximations have different
signs, which we noticed first in the tables.

\subsection{Simpson's Rule}
\index{Simpson's rule}

Our goal is a calculation scheme for integrals 
%
\mar{Combine good approximations\ldots}
%
that gives an error is as small as possible.  The trapezoidal and the
midpoint approximations are both good---but we can combine them to get
something even better.  Here is why.  The tables and the figure above
indicate that the {\em errors\/} in the two approximations have
opposite signs, and the midpoint error is only about half the size of
the trapezoidal error (in absolute value).  Thus, if we form the sum
	$$
2 \times \mbox{midpoint approximation} 
   + \mbox{trapezoidal approximation}, 
	$$
then most of the error will cancel.  
%
\mar{\ldots so that most of the error cancels}
%
Now this sum is approximately three times the value of the integral,
because each term in it approximates the integral itself.  Therefore,
if we divide by three, then
	$$
	\tfrac{2}{3} \times \mbox{midpoint approximation} + 
	\tfrac{1}{3} \times \mbox{trapezoidal approximation}
	$$
should be a superb approximation to the integral.  

	Let's try this approximation on our test integral
	$$
	\int_1^3 \dfrac{1}{x}\,dx
	$$
with $n = 100$ subintervals.  Using the numbers from the table on
page~\pageref{trap}, we obtain
	$$
	\tfrac{2}{3} \times 1.09859747 + 
	\tfrac{1}{3} \times 1.09864191 = 1.098612283,
	$$ 
which gives the value of the integral accurate to 8 decimal places.
This method of approximating integrals is called {\bf Simpson's
  rule}.\index{Simpson's rule}

We can use the program RIEMANN\index{program!RIEMANN} to do the calculation.
%
\mar{Using Riemann sums to carry out\\ Simpson's rule}	
%
First calculate left and right Riemann sums, and take their average.
That is the trapezoidal approximation.  Then calculate the midpoint
Riemann sum.  Since
	$$
	\mbox{trapezoid} = \tfrac{1}{2} \times \mbox{left} + 
		\tfrac{1}{2} \times \mbox{right},
	$$
Simpson's rule\index{Simpson's rule} reduces to this combination of
left, right, and midpoint sums:
\begin{align*} 
\lefteqn{\tfrac{2}{3} \times \mbox{midpoint} + 
	\tfrac{1}{3} \times \mbox{trapezoidal}} \qquad\quad  \\
  &=  \tfrac{2}{3} \mbox{midpoint} + \tfrac{1}{3}  
		\left( \tfrac{1}{2} \times \mbox{left} + 
		\tfrac{1}{2} \times \mbox{right} \right) \\[3pt]
  &= \tfrac{2}{3} \times \mbox{midpoint} + 
		\tfrac{1}{6} \times \mbox{left} + 
		\tfrac{1}{6} \times \mbox{right} \\[3pt]
  &= \tfrac{1}{6} \left( 4 \times \mbox{midpoint} + 
		\mbox{left} + \mbox{right} 
		\vphantom{\tfrac{1}{6}}\right) 
\end{align*}
\begin{center}
\begin{picture}(350,60)
\put(0,0){\framebox(350,60)
{
\begin{minipage}{340pt}			
\begin{center} 
{\bf Simpson's rule}: \\ [3pt]
	$\dint f(x) \, dx \approx \dfrac{1}{6}
	\left( \vphantom{\tfrac{1}{6}}\mbox{left sum} + 
	\mbox{right sum} + 4 \times \mbox{midpoint sum} \right)$
\end{center}
\end{minipage}
}}
\end{picture}
\end{center}

You can get even more accuracy 
%
\mar{The accuracy of Simpson's rule}	
%
if you keep track of more digits in the left, right, and midpoint
Riemann sums.  For example, if you estimate
	$$
	\int_1^3 \dfrac{1}{x}\,dx = \ln{3} = 1.098\,612\,288\,668\ldots 
	$$
to 14 decimal places, you will get the following.
	\begin{center}
	\begin{tabular}{rl}
	left: 		& 1.105\,308\,583\,647\,79 \\
	right: 		& 1.091\,975\,250\,314\,45 \\
	midpoint:	& 1.098\,597\,475\,005\,31
	\end{tabular}
	\end{center}
When combined these give the estimate 1.098\,612\,288\,997\,3, which
differs from the true value by less than $3.3 \times 10^{-10}$.  In
other words, the calculation is actually correct to 9 decimal places.

It is possible to get a bound on the error 
%
\mar{An error bound for Simpson's rule}\index{error bound}
%
produced by using Simpson's rule to estimate the value of
	$$
	\int_a^b f(x) \, dx.
	$$
(See the discussion of error bounds in chapter~6.3.)  Specifically,
	$$
	\left| \, \int_a^b f(x) \, dx - \mbox{Simpson's rule } \right| \leq
		\dfrac{M (b - a)^5}{2880 \, n^4},
	$$ 
where $n$ is the number of subintervals used in the Riemann sums and $M$ is a bound on the size of the fourth derivative of $f$:
	$$
	|f^{(4)}(x)| \leq M \qquad \mbox{for all $a \leq x \leq b$.}
	$$

\newpage

The most important factor in the error bound 
%
\mar{The crucial factor $n^{-4}$}
%
is the $n^4$ that appears in the denominator.  In our example $n = 100
= 10^2$, and this leads to the factor $1/n^4 = 10^{-8}$ in the error
bound.  As we saw, the actual error was less than $10^{-9}$.
Essentially, $n$ is the number of computations we do, and error bound
tells how many decimal places of accuracy we can count on.  According
to the error bound, a ten-fold increase in the number of computations
produces four more decimal places of accuracy.  {\em That\/} is why
Simpson's method is efficient.

\subsection{Exercises}

\vspace{-10pt}

Because Simpson's rule is so efficient, we can use it to get
accurate values of some of the fundamental constants of mathematics.
For example, since
	$$
	4 \cdot \int_0^1 \dfrac{dx}{1 + x^2} =  
		\left. 4\arctan(x) \ \vphantom{\dfrac{1}{2}} \right|_0^1 = 
		4\arctan(1) = 4 \cdot \dfrac{\pi}{4} = \pi,
	$$
we can estimate the value of $\pi$ by using Simpson's rule to
approximate this integral.

\num Evaluate the expression above (including the factor of 4) use
Simpson's rule with $n = 2$, 4, 8, and 16.  How accurate is each of
these estimates of $\pi$; that is, how many decimal places of each
estimate agree with the true value of $\pi$?

\numa   Over the interval $0 \leq x \leq 1$ it is true that
$$
|f^{(4)}(x)| \leq 96 \qquad \mbox{when} \qquad f(x) = \dfrac{4}{1+x^2}.
$$
(You don't need to show this, but how might you do it?)  Use this
bound to show that $n = 256 = 2^8$ will guarantee that you can find
the first 10 decimals of $\pi$ by using the method of the previous
question.

\sub Show that if $n = 128 = 2^7$ then the error bound for Simpson's
rule does {\em not\/} guarantee that you can find the first 10
decimals of $\pi$ by the same method.

\sub Run Simpson's rule with $n = 2^7$ to estimate $\pi$.  How many
decimal places {\em are\/} correct?  Does this surprise you?  In fact
the error bound is too timid: it says that the error is no larger than
the bound it gives, but the actual error may be much smaller.  From
your work in part (a), which power of 2 is sufficient to get 10
decimal places accuracy?

\numa In chapter 6.3 (page~\pageref{int error}) a left Riemann sum for		
$$
\int_0^1 e^{-x^2} \, dx
$$ 
with 1000 equal subdivisions gave 3 decimal places accuracy.  How many
subdivisions $n$ are needed to get that much accuracy using Simpson's
rule?  Let $n$ be a power of 2.  Start with $n = 1$ and increase $n$
until three digits stabilize.

\sub If you use Simpson's rule with $n = 1000$ to estimate this
integral, how many digits stabilize?

\num On page~\pageref{midpoint for x-cube} a midpoint Riemann sum with
$n = 10000$ shows that
$$
\int_1^3 \sqrt{1 + x^3} \, dx = 6.229959\ldots .
$$
How many subdivisions $n$ are needed to get this much accuracy using
Simpson's rule?  Start with $n = 1$ and keep doubling it until seven
digits stabilize.

%%%%%%%%%%%%%%%%%%%%%%%%%%%%%%%%%%%%%%%%%%%%%%%%%%%%%%%%%%%%%%%%%%%%%%%%
%                                                                      %
%                               Section 7                              %
%                                                                      %
%%%%%%%%%%%%%%%%%%%%%%%%%%%%%%%%%%%%%%%%%%%%%%%%%%%%%%%%%%%%%%%%%%%%%%%%

\section{Improper Integrals}\index{improper integral}\index{integral!improper}

\subsection{The Lifetime of Light Bulbs}\index{light bulbs!lifetime of}

Ordinary light bulbs are supposed to burn about 700 hours, 
%
\mar{The lifetime of a light bulb is unpredictable}	
%
but of course some last longer while others burn out more quickly.  It
is impossible to know, in advance, the lifetime of a particular bulb
you might buy, but it is possible to describe what happens to a large
batch of bulbs.

Suppose we take a batch of 1000 light bulbs, start them burning at the
same time, and note how long it takes each one to burn out.  Let
	$$
	L(t) = \mbox{fraction of bulbs that burn out before 
		$t$ hours}
	$$
Then $L(t)$ might have a graph that looks like this:

\begin{picture}(340,160)(15,-10)
	\put(50,20){\vector(1,0){255}}
	\put(50,20){\vector(0,1){108}}
	\put(300,25){$t$}
	\put(56,122){$L$}
	\put(90,20){\line(0,1){3}}
	\put(130,20){\line(0,1){3}}
	\put(170,20){\line(0,1){3}}
	\put(210,20){\line(0,1){3}}
	\put(250,20){\line(0,1){3}}
	\put(50,8){\makebox[0pt]{0}}
	\put(90,8){\makebox[0pt]{250}}
	\put(130,8){\makebox[0pt]{500}}
	\put(170,8){\makebox[0pt]{750}}
	\put(210,8){\makebox[0pt]{1000}}
	\put(250,8){\makebox[0pt]{1250}}
	\put(290,8){hours}
	\put(50,37){\line(1,0){3}}
	\put(50,54){\line(1,0){3}}
	\put(50,71){\line(1,0){3}}
	\put(50,88){\line(1,0){3}}
	\put(50,105){\line(1,0){3}}
	\put(46,33){\makebox[0pt][r]{.2}}
	\put(46,51){\makebox[0pt][r]{.4}}
	\put(46,68){\makebox[0pt][r]{.6}}
	\put(46,85){\makebox[0pt][r]{.8}}
	\put(46,102){\makebox[0pt][r]{1.0}}
	\put(46,122){\makebox[0pt][r]{fraction}}
	\begin{thicklines}
	\bezier{90}(50,20)(90,32)(120,62)
	\bezier{200}(120,62)(160,103)(300,103)
	\end{thicklines}
	\multiput(56,105)(4,0){62}{\circle*{1}}
\end{picture}
In this example, $L(400) \approx .5$, so about half the bulbs burned out before 400 hours.  Furthermore, all but a few have burned out by 1250 hours.  

Manufacturers are very concerned 
%
\mar{The burnout rate}	
%
about the way the lifetime of light bulbs varies.  They study the
output of their factories on a regular basis.  It is more common,
though, for them to talk about the {\em rate\/} $r$ at which bulbs
burn out.  The rate varies over time, too.  In fact, in terms of $L$,
$r$ is just the derivative
	$$
	r(t) = L'(t) \quad \mbox{bulbs per hour.}
	$$
However, if we {\em start\/} with the rate $r$, then we get $L$ as the integral
	$$
	L(t) = \int_0^t r(s) \, ds.
	$$
This is yet another consequence of 
%
\mar{Lifetime is the integral of burnout rate\ldots}
%
the fundamental theorem of calculus\index{fundamental theorem of
  calculus}.  The integral expression is quite handy.  For example,
the fraction of bulbs that burn out between $t= a$ hours and $t = b$
hours is
	$$
	L(b) - L(a) = \int_a^b r(s) \, ds.
	$$

We can even use the integral to say that all the bulbs burn out eventually:
	$$
	L(t) = \int_0^t r(s) \, ds = 1 \qquad 
		\mbox{when $t$ is sufficiently large.}
	$$
In practice $r$ is the average burnout rate 
%
\mar{\ldots but there is no upper limit to the lifetime}
%
for many batches of light bulbs, so we can't identify the precise
moment when $L$ becomes 1.  All we can really say is
\[
	L(\infty) = \int_0^{\infty} r(s) \, ds.
\]
This is called an {\bf improper integral}\index{improper integral}\index{integral!improper}, 
%
\mar{An integral is {\em improper\/} if the domain of integration is
  infinite}
%
because it cannot be calculated directly: its ``domain of
integration'' is infinite.  By definition, its value is obtained as a
limit of ordinary integrals:
	$$
	\int_0^{\infty} r(s) \, ds = 
		\lim_{b \to \infty} \int_0^b r(s) \, ds.
	$$

The \label{`normal'} {\bf normal density function}\index{function!normal density} of probability theory provides us with another
example of an improper integral.  In a simple form, the function
itself is
$$
f(x) = \dfrac{1}{\sqrt{2\pi}} \exp\left(-\dfrac{x^2}{2}\right).
$$

(Recall, $\exp(x) = e^x$.)  If $x$ is any {\em normally distributed
  quantity\/} whose average value is 0, then the probability that a
randomly chosen value of $x$ lies between the numbers $a$ and $b$ is
	$$
 	\int_a^b f(x) \, dx.
	$$
Since the probability that $x$ lies {\em somewhere\/} 
%
\mar{The normal probability distribution involves\\ an improper integral}
%
on the $x$-axis is 1, we have
	$$
 	\int_{-\infty}^{\infty} f(x) \, dx = 1.
	$$
This is an improper integral\index{improper integral}, and its value
is defined by the limit
	$$
 	\int_{-\infty}^{\infty} f(x) \, dx = 
		\lim_{b \to \infty} \int_{-b}^b f(x) \, dx.
	$$
In the exercises you will have a chance to evaluate this integral.


\subsection{Evaluating Improper Integrals}\index{improper integral!evaluating}\index{integral!improper!evaluating}

An integral with an infinite domain of integration is 
%
\mar{An integral is also improper if its integrand becomes infinite}	
%
only one kind of improper integral.  A second kind has a finite domain
of integration, but the integrand\index{integrand} becomes infinite on
that domain.  For example,
	$$
	\int_0^1 \dfrac{dx}{x} \qquad \mbox{and} \qquad
		\int_0^1 \ln{x} \, dx 
	$$
are both improper in this sense.  In both cases, the integrand becomes
infinite as $x \to 0$.  Because the difficulty lies at the endpoint 0,
we define
	$$
	\int_0^1 \dfrac{dx}{x} = 
		\lim_{a \to 0} \int_a^1 \dfrac{dx}{x}.
	$$
More generally, 
	$$
	\int_a^b f(x) \, dx 
	$$
is an improper integral if $f(x)$ becomes infinite at some point $c$
in the interval $[a, b]$.  In that case we define
	$$
	\int_a^b f(x) \, dx = \lim_{q \to 0} 
		\left(\int_a^{c-q} f(x) \, dx + 
			\int_{c+q}^b f(x) \, dx\right).
	$$
In effect, we avoid the bad spot but ``creep up'' on it in the limit.

Indefinite integrals\index{indefinite integral}---that is, 
%
\mar{Antiderivatives help find improper integrals}	
%
antiderivatives\index{antiderivative}---can be a great help in
evaluating improper integrals.  Here are some examples.
\bigskip

\noindent
{\bf Example 1}.  We can evaluate 
	$$
	\int_0^{\infty} e^{-x} \, dx 
	$$
by noting first that $\dint e^{-x} \, dx = - e^{-x}$.  Therefore
	$$
	\int_0^b e^{-x} \, dx = \left. - e^{-x} \right|_0^b =
		-e^{-b} - \left(-e^{-0}\right) = 1 - e^{-b} 
	$$
and
	$$
	\int_0^{\infty} e^{-x} \, dx = 
		\lim_{b \to \infty} \int_0^b e^{-x} \, dx = 
		\lim_{b \to \infty} \left(1-e^{-b}\right) = 1.
	$$

\medskip

\noindent
{\bf Example 2}.  To evaluate $\dint_0^1 \ln{x} \, dx$, we use the
indefinite integral
	$$
	\int \ln{x} \, dx = x\, \ln{x} - x.
	$$
Thus
	$$
	\int_a^1 \ln{x} \, dx = 
		\left. x\,\ln{x} - x
			\rule{0pt}{14pt}\,\right|_a^1 = 
		-1 - \left( a\ln{a} - a \right) = 
		a-1 - a\ln{a}.
	$$
By direct calculation (using a graphing package, for instance) we can find 
	$$
	\lim_{a \to 0} \ a\ln{a} = 0;
	$$
therefore
\[
\int_0^1 \ln{x} \, dx 
  = \lim_{a \to 0} \int_a^1 \ln{x} \, dx
  = \lim_{a \to 0} \ a-1 - a\ln{a}
  = -1.
\]

You should not assume that an improper integral\index{improper
  integral!finite value of} always has a finite value, though.
Consider the next example.

\bigskip

\noindent
{\bf Example 3}. $\displaystyle \left. \dint_0^1 \dfrac{dx}{x} 
= \lim_{a \to 0} \int_a^1 \dfrac{dx}{x}
= \lim_{a \to 0} \ln(x) \,\right|_a^1
= \lim_{a \to 0} \left( \ln(1)-\ln(a) \right)
= \infty.
$


\bigskip

\noindent
This is forced because ${\displaystyle \lim_{a \to 0} \ \ln(a) =
  -\infty}$, which you can see from the graph of the logarithm
function.

\subsection{Exercises}

\vspace{-10pt} 

\num Find the value of each of the following improper
integrals.  (The value may be $\infty$.)

\noindent
\begin{minipage}{150pt}
\sub $\dint_{-\infty}^0 e^x \, dx$

\sub $\dint_1^{\infty} \dfrac{du}{u}$


\sub $\dint_0^1 \dfrac{dy}{y^2}$

\sub $\dint_0^{\pi/2} \tan{x} \, dx$
\end{minipage}
\hfill
\begin{minipage}{150pt}
	
\sub $\dint_0^{\infty} x e^{-x} \, dx$
	
\sub $\dint_1^{\infty} \dfrac{du}{u^2}$
	
\sub $\dint_0^{\infty} \dfrac{x}{1 + x^2} \, dx$
	 
\sub $\dint_1^3 \dfrac{x}{x^2 - 1}\, dx$
\end{minipage}
\hfill
\mbox{}

\num Use the reduction formula for $\dint \dfrac{dx}{(1+x^2)^n}$ you
found on page~\pageref{1plusx2} to find the exact value of
$\dint_0^{\infty} \dfrac{dx}{(1+x^2)^{10}}$

\subsubsection{The normal density function}

The next two questions concern the improper integral
	$$
	\dfrac{1}{\sqrt{2\pi}} \int_{-\infty}^{\infty}
		e^{-x^2/2} \, dx 
	$$
of the normal density function\index{function, normal density} defined on page~\pageref{`normal'}.  The goal is to determine the value of this integral.

\num First, use RIEMANN\index{program!RIEMANN} to estimate
the value of
	$$
	\dfrac{1}{\sqrt{2\pi}} \int_{-b}^{b}
		e^{-x^2/2} \, dx 
	$$
when $b$ has the different values 1, 10, 100, and 1000.  On the basis
of these results, estimate
	$$
	\lim_{b \to \infty} \dfrac{1}{\sqrt{2\pi}} \int_{-b}^{b}
		e^{-x^2/2} \, dx.
	$$
This gives one estimate of the value of the improper integral.

\numa To construct a second estimate, begin by sketching the graph of
the normal density function
	$$
	f(x) = \dfrac{1}{\sqrt{2\pi}} e^{-x^2/2} 
	$$
on an interval centered at the origin.  Use the graph to argue that 
	$$
	\dfrac{1}{\sqrt{2\pi}} \int_{-b}^{b}
		e^{-x^2/2} \, dx =
	2 \left( \dfrac{1}{\sqrt{2\pi}} \int_{0}^{b}
		e^{-x^2/2} \, dx \right) =
	\sqrt{2/\pi} \int_{0}^{b}
		e^{-x^2/2} \, dx 
	$$
and therefore
	$$
	\dfrac{1}{\sqrt{2\pi}} \int_{-\infty}^{\infty}
		e^{-x^2/2} \, dx =
	\sqrt{2/\pi} \int_{0}^{\infty}
		e^{-x^2/2} \, dx.
	$$

\sub  Now consider the accumulation function
$$
F(t) = \sqrt{2/\pi} \int_{0}^{t} e^{-x^2/2} \, dx.
$$
We want to find $F(\infty) = {\displaystyle \lim_{t \to \infty}
  F(t)}$.  According to the fundamental theorem of calculus, $y =
F(t)$ satisfies the initial value problem
	$$
	\dfrac{dy}{dt} = \sqrt{2/\pi} \cdot
		e^{-t^2/2}, \qquad
		y(0) = 0.
	$$
Use a differential equation solver (e.g., PLOT\index{program!PLOT})
to graph the solution $y = F(t)$ to this problem.  From the graph
determine
	$$
	F(\infty) = \lim_{t \to \infty} F(t).
	$$

\sub Does your results in part (b) and question 2 agree?  Do they
agree with the value the text claims for the improper integral.
(Remember, the value is the probability that a randomly chosen number
will lie {\em somewhere\/} on the number line between $-\infty$ and
$+\infty$.)


\subsubsection{The gamma function}\index{function, gamma}

The {\bf factorial function}\index{function!factorial} is defined for
a positive integer $n$ by the formula
	$$
	n! = n \cdot (n-1) \cdot (n-2) \cdot \cdots \cdot 
		3 \cdot 2 \cdot 1.
	$$ 
For example, $1! = 1$, $2! = 2$, $3! = 6$, $4! = 24$, and $10! =
3628800$.  The factorial function is used often in diverse
mathematical contexts, but its use is sometimes limited by the fact
that it is defined only for positive integers.  How might the function
be defined on an expanded domain, so that we could deal with
expressions like $tfrac{1}{2}!$, for example?  The {\em gamma
  function\/} answers this question.
\bigskip

The {\bf gamma function}\index{gamma function} $\Gamma(x)$ is defined
by the improper integral
$$
\Gamma(x) = \int_0^{\infty} e^{-t} t^{x - 1} \, dt.
$$ Notice that $t$ is the active variable in this integral.  While the
integration is being performed, $x$ is treated as a constant; for the
integral to converge, we need $x > 0$.

\num  Show that $\Gamma(1) = 1$.

\num  Using integration by parts, show that 
$\Gamma(x + 1) = x \cdot \Gamma(x)$.  You may use the fact that
	$$
	\dfrac{t^p}{e^t} \to 0 \ \mbox{as} \ t \to +\infty.
	$$
\smallskip

\noindent
The property $\Gamma(x + 1) = x \cdot \Gamma(x)$ makes the gamma
function like the factorial function, because
	$$
	(n+1)! = (n+1) \cdot 
		\underbrace{n \cdot (n-1) \cdot (n-2) \cdot 
			\cdots \cdot 3 \cdot 2 \cdot 1}_
		{\displaystyle = n!}
		= (n+1) \cdot n!.
	$$
Notice there is a slight difference, though.  We explore this now.

\num Using the property $\Gamma(x + 1) = x \cdot \Gamma(x)$, calculate
$\Gamma(2)$, $\Gamma(3)$, $\Gamma(4)$, $\Gamma(5)$, and $\Gamma(6)$.
On the basis of this evidence, fill in the blank:
	$$
	\mbox{For a positive integer $n$,} \qquad
		\Gamma(n) = \rule[-2pt]{30pt}{.3pt}\, ! 
	$$
Using this relation, give a computable meaning to the
expression~$\tfrac{1}{2}!$.


\num Estimate the value of $\Gamma(1/2)$.  (Exercises 2 and 3 above
offer two ways to estimate the value of an improper integral.)

\num In fact, $\Gamma(1/2) = \sqrt{\pi}$ exactly.  You can show this
by employing several of the techniques developed in this chapter.
Start with
	$$
	\Gamma(1/2) = 
		\int_0^{\infty} e^{-t} t^{-\frac{1}{2}} \, dt.
	$$

\sub Make the substitution $u = (2t)^{\tfrac{1}{2}}$ and show that the
integral becomes
	$$
	\Gamma(1/2) = 
		\sqrt{2} \int_0^{\infty} e^{-u^2/2} du.
	$$

\sub  From exercise 4 you know 
	$$
	\sqrt{2/\pi} \int_{0}^{\infty}
		e^{-u^2/2} \, du = 1.
	$$
(Check this.)  Now, using some algebra, show $\Gamma(1/2) = \sqrt{\pi}$.

\sub Compare your estimate for $\Gamma(1/2)$ from exercise 7 with the
exact value $\sqrt{\pi}$.

\numa Determine the exact values of $\Gamma(3/2)$ and $\Gamma(5/2)$.

\sub In exercise 6 you gave a meaning to the expression
$\tfrac{1}{2}!$; can you now give it an exact value?

\newpage
%%%%%%%%%%%%%%%%%%%%%%%%%%%%%%%%%%%%%%%%%%%%%%%%%%%%%%%%%%%%%%%%%%%%%%%%
%                                                                      %
%                               Section 8                              %
%                                                                      %
%%%%%%%%%%%%%%%%%%%%%%%%%%%%%%%%%%%%%%%%%%%%%%%%%%%%%%%%%%%%%%%%%%%%%%%%
\section{Chapter Summary}

\subsection{The Main Ideas}


\begin{itemize}

\item A function $F$ is an {\bf antiderivative} of $f$ if $F'=f$.
  {\em Every\/} antiderivative of $f$ is equal to $F+C$ for some
  appropriately chosen constant $C$.  We write $\int \! f(x) \, dx =
  F(x) +C$.

\item Differentiation rules for combinations of functions yield
  corresponding anti-differentiation rules.  Among these are the {\bf
    constant multiple} and {\bf addition} rules.  The chain rule for
  differentiation corresponds to {\bf integration by substitution}.
  The product rule for differentiation corresponds to {\bf integration
    by parts}.


\item The derivative of a function and of its inverse are reciprocals.
  When $y=f(t)$ and $t=g(y)$ are inverses:
$$
\dfrac{dt}{dy} = \frac{1}{dy/dt}.
$$

\item In some cases, the method of {\bf separation of variables} can
  be used to find a {\em formula\/} for the solution of a differential
  equation.

\item A numerical method for estimating an integral is {\bf efficient}
  if it gets accurate results with relatively few calculations.  The
  {\bf trapezoidal approximation} is the average of a left and a right
  {\em endpoint} Riemann sum and is more efficient than either.  {\em
    Midpoint} Riemann sums are even more efficient than trapezoidal
  approximations.

\item The most efficient method developed in this chapter is {\bf
  Simpson's rule}.  Simpson's rule approximates an integral by
$$
\int_a^b f(x) \, dx \approx \frac{1}{6}\left( \mbox{left sum} + \mbox{right sum} + 4 \times \mbox{midpoint sum} \right).
$$

\item An {\bf improper integral} is one that cannot be calculated
  directly.  The problem may be that its ``domain of integration'' is
  infinite or that the integrand becomes infinite on that domain.  Its
  value is obtained as a limit of ordinary integrals.

\end{itemize}


\subsection{Expectations}

\begin{itemize}

\item  You should be able to find antiderivatives of basic functions.

\item You should be able to find antiderivatives of combinations of
  functions using the {\bf constant multiple} and {\bf addition}
  rules, as well as the {\bf method of substitution} and {\bf
    integration by parts}.

\item You should be able to rewrite an integrand given as a quotient
  using the method of {\bf partial fractions}.

\item You should be able to express the derivative of an invertible
  function in terms of the derivative of its inverse.  In particular,
  you should be able to differentiate the {\bf arctangent}, {\bf
    arcsine} and {\bf arccosine} functions.

\item You should be able to solve a differential equation using the
  method of {\bf separation of variables}.

\item You should be able to adapt the program RIEMANN to approximate
  integrals using the {\bf trapezoid rule} and {\bf Simpson's rule}.

\item You should be able to find the value of an {\bf improper
  integral} as the limit of ordinary integrals.

\end{itemize}

